% claim.tex -- AI4Math claim of record (AI-generated, 2026-09-30; instantiated
% from harness/templates/claim/claim.tex per SOP 08b).
% Machine-readable theorem anchor is the "% LEAN:" line; scripts/paper-lint.sh
% and the claim audit check it. All Lean listings are verbatim, byte-identical
% contiguous excerpts of frozen artifacts of tasks/20260930-comb12-grid-indep:
%   statement  : formalized/statement.lean (61 content lines, LF, sha256 b726f0d8...;
%                the theorem-head excerpt stops before the frozen file's sorry
%                placeholder and native_decide data-check lines)
%   interface  : formalized/interface.lean (8042 bytes, sha256 66c96c3e...);
%                excerpt = lines 18-163 (frozen segment + attack interface +
%                matrix certificates; the roster/attack-notes docstrings that
%                mention banned words in warning text are excluded)
%   proof file : attempts/final-c_z_rec.lean (1188 lines + trailing newline,
%                54264 bytes, mixed
%                line endings on the working disk: lines 1-485 and 953-1188
%                LF (content lines; the split-chunk count 1189 includes the
%                trailing newline), 486-952 CRLF; spliced LF-normalised per
%                the chorded/ladder precedent, byte-exact original shipped
%                in proofs/)
% Splice/verification: audit/inspect-tex.py (--splice, then --verify;
% record in audit/listing-verify.txt).
\documentclass[11pt]{article}

\usepackage[T1]{fontenc}
\usepackage[utf8]{inputenc}
\usepackage{amsmath,amssymb,amsthm}
\usepackage{graphicx}
\usepackage{hyperref}
\usepackage{xurl}
\setlength{\emergencystretch}{3em}
% lstlean.tex — Lean style for the listings package.
% Lineage: Jeremy Avigad (2015), adapted from Assia Mahboubi's SSR style;
% inspired by Olivier Verdier's unixode.sty for the Unicode literate table.
% No CTAN package or upstream repo exists; papers carry this file in their
% source package (cap set arXiv:1907.01449, sphere eversion arXiv:2210.07746).
% AI4Math adaptation (AI-generated, 2026-09-26): sorry/admit/sorryAx elevated
% to a dedicated error tier, matching the pipeline's 0-sorry constitution.
\usepackage{listings}
\usepackage{xcolor}

\definecolor{keywordcolor}{rgb}{0.7, 0.1, 0.1}   % red keywords
\definecolor{tacticcolor}{rgb}{0.0, 0.1, 0.6}    % blue tactics
\definecolor{commentcolor}{rgb}{0.4, 0.4, 0.4}   % grey comments
\definecolor{symbolcolor}{rgb}{0.0, 0.1, 0.6}    % blue symbols
\definecolor{sortcolor}{rgb}{0.1, 0.5, 0.1}      % green sorts
\definecolor{errorcolor}{rgb}{0.6, 0, 0}         % dark red: sorry/admit
\definecolor{stringcolor}{rgb}{0.5, 0.1, 0.5}    % purple strings

\lstdefinelanguage{lean}{
  % Anything between $ becomes LaTeX math mode
  mathescape=false,
  % Comments may or not include Latex commands
  texcl=false,
  % keywords, list taken from lean-syntax.txt
  morekeywords=[1]{import, prelude, protected, private, noncomputable,
    definition, meta, renaming, hiding, parameter, parameters, begin,
    constant, constants, lemma, variable, variables, theory, print,
    theorem, example, open, section, namespace, universe, universes, end,
    modifier, modifiers, using, export, exposing, axiom, inductive,
    coinductive, with, structure, class, instance, attribute, local, where,
    by, have, show, let, in, if, then, else, match, do, for, fun, assume,
    from, take, calc, include, omit, infix, infixl, infixr, notation,
    macro, syntax, elab, set_option, initialize, deriving, opaque, abbrev,
    mutual, partial, scoped, termination_by, decreasing_by},
  % Sorts
  morekeywords=[2]{Sort, Type, Prop},
  % tactics, list taken from lean-syntax.txt (tactic keywords)
  morekeywords=[3]{intro, intros, apply, exact, refine, rfl, simp, simp_all,
    simpa, simp_rw, rw, rewrite, erewrite, ring, ring_nf, omega, decide,
    norm_num, norm_cast, induction, cases, rcases, obtain, constructor,
    ext, funext, conv, congr, field_simp, linarith, nlinarith, positivity,
    tauto, trivial, aesop, use, by_cases, by_contra, push_neg, contrapose,
    symm, trans, assumption, contradiction, exfalso, specialize,
    generalize, revert, clear, rename_i, subst, unfold, delta, change,
    convert, split_ifs, interval_cases, fin_cases, zify, gcongr,
    nth_rewrite, suffices, generalizing, at, only, native_decide},
  % warnings - constitution tier: must never survive into a final proof
  morekeywords=[4]{sorry, admit, sorryAx},
  literate=
    {α}{{\ensuremath{\alpha}}}1
    {β}{{\ensuremath{\beta}}}1
    {γ}{{\ensuremath{\gamma}}}1
    {δ}{{\ensuremath{\delta}}}1
    {ε}{{\ensuremath{\varepsilon}}}1
    {ζ}{{\ensuremath{\zeta}}}1
    {η}{{\ensuremath{\eta}}}1
    {θ}{{\ensuremath{\theta}}}1
    {ι}{{\ensuremath{\iota}}}1
    {κ}{{\ensuremath{\kappa}}}1
    {λ}{{\ensuremath{\lambda}}}1
    {μ}{{\ensuremath{\mu}}}1
    {ν}{{\ensuremath{\nu}}}1
    {ξ}{{\ensuremath{\xi}}}1
    {π}{{\ensuremath{\pi}}}1
    {ρ}{{\ensuremath{\rho}}}1
    {σ}{{\ensuremath{\sigma}}}1
    {τ}{{\ensuremath{\tau}}}1
    {υ}{{\ensuremath{\upsilon}}}1
    {φ}{{\ensuremath{\varphi}}}1
    {χ}{{\ensuremath{\chi}}}1
    {ψ}{{\ensuremath{\psi}}}1
    {ω}{{\ensuremath{\omega}}}1
    {Γ}{{\ensuremath{\Gamma}}}1
    {Δ}{{\ensuremath{\Delta}}}1
    {Θ}{{\ensuremath{\Theta}}}1
    {Λ}{{\ensuremath{\Lambda}}}1
    {Ξ}{{\ensuremath{\Xi}}}1
    {Π}{{\ensuremath{\Pi}}}1
    {Σ}{{\ensuremath{\Sigma}}}1
    {Φ}{{\ensuremath{\Phi}}}1
    {Ψ}{{\ensuremath{\Psi}}}1
    {Ω}{{\ensuremath{\Omega}}}1
    {ℵ}{{\ensuremath{\aleph}}}1
    {∀}{{\ensuremath{\forall}}}1
    {∃!}{{\ensuremath{\exists}!}}1
    {∃}{{\ensuremath{\exists}}}1
    {∈}{{\ensuremath{\in}}}1
    {∉}{{\ensuremath{\notin}}}1
    {∘}{{\ensuremath{\circ}}}1
    {∩}{{\ensuremath{\cap}}}1
    {∪}{{\ensuremath{\cup}}}1
    {⊆}{{\ensuremath{\subseteq}}}1
    {⊂}{{\ensuremath{\subset}}}1
    {⊇}{{\ensuremath{\supseteq}}}1
    {⊃}{{\ensuremath{\supset}}}1
    {⊄}{{\ensuremath{\nsubseteq}}}1
    {∅}{{\ensuremath{\emptyset}}}1
    {∧}{{\ensuremath{\wedge}}}1
    {∨}{{\ensuremath{\vee}}}1
    {¬}{{\ensuremath{\neg}}}1
    {⊤}{{\ensuremath{\top}}}1
    {⊥}{{\ensuremath{\bot}}}1
    {↔}{{\ensuremath{\leftrightarrow}}}1
    {→}{{\ensuremath{\rightarrow}}}1
    {←}{{\ensuremath{\leftarrow}}}1
    {↑}{{\ensuremath{\uparrow}}}1
    {⇒}{{\ensuremath{\Rightarrow}}}1
    {⇐}{{\ensuremath{\Leftarrow}}}1
    {⟹}{{\ensuremath{\Longrightarrow}}}1
    {↦}{{\ensuremath{\mapsto}}}1
    {≤}{{\ensuremath{\leq}}}1
    {≥}{{\ensuremath{\geq}}}1
    {≠}{{\ensuremath{\neq}}}1
    {≈}{{\ensuremath{\approx}}}1
    {≡}{{\ensuremath{\equiv}}}1
    {≃}{{\ensuremath{\simeq}}}1
    {≅}{{\ensuremath{\cong}}}1
    {×}{{\ensuremath{\times}}}1
    {·}{{\ensuremath{\cdot}}}1
    {±}{{\ensuremath{\pm}}}1
    {⊕}{{\ensuremath{\oplus}}}1
    {⊗}{{\ensuremath{\otimes}}}1
    {⋆}{{\ensuremath{\star}}}1
    {▸}{{\ensuremath{\blacktriangleright}}}1
    {⟨}{{\ensuremath{\langle}}}1
    {⟩}{{\ensuremath{\rangle}}}1
    {⦃}{{\ensuremath{\{\!\mid}}}1
    {⦄}{{\ensuremath{\mid\!\}}}}1
    {⟦}{{\ensuremath{[\![}}}1
    {⟧}{{\ensuremath{]\!]}}}1
    {⌊}{{\ensuremath{\lfloor}}}1
    {⌋}{{\ensuremath{\rfloor}}}1
    {⌈}{{\ensuremath{\lceil}}}1
    {⌉}{{\ensuremath{\rceil}}}1
    {√}{{\ensuremath{\surd}}}1
    {∣}{{\ensuremath{\mid}}}1
    {∤}{{\ensuremath{\nmid}}}1
    {∎}{{\ensuremath{\blacksquare}}}1
    {⋯}{{\ensuremath{\cdots}}}1
    {…}{{\ensuremath{\ldots}}}1
    {∑}{{\ensuremath{\sum}}}1
    {∏}{{\ensuremath{\prod}}}1
    {⁻¹}{{\ensuremath{^{-1}}}}1
    {¹}{{\ensuremath{^1}}}1
    {²}{{\ensuremath{^2}}}1
    {³}{{\ensuremath{^3}}}1
    {ⁿ}{{\ensuremath{^n}}}1
    {ᵐ}{{\ensuremath{^m}}}1
    {₀}{{\ensuremath{_0}}}1
    {₁}{{\ensuremath{_1}}}1
    {₂}{{\ensuremath{_2}}}1
    {₃}{{\ensuremath{_3}}}1
    {₄}{{\ensuremath{_4}}}1
    {₅}{{\ensuremath{_5}}}1
    {ₙ}{{\ensuremath{_n}}}1
    {ₘ}{{\ensuremath{_m}}}1
    {ₖ}{{\ensuremath{_k}}}1
    {ᵢ}{{\ensuremath{_i}}}1
    {ⱼ}{{\ensuremath{_j}}}1
    {ℤ}{{\ensuremath{\mathbb{Z}}}}1
    {ℕ}{{\ensuremath{\mathbb{N}}}}1
    {ℚ}{{\ensuremath{\mathbb{Q}}}}1
    {ℝ}{{\ensuremath{\mathbb{R}}}}1
    {ℂ}{{\ensuremath{\mathbb{C}}}}1
    {𝔽}{{\ensuremath{\mathbb{F}}}}1
    {𝒫}{{\ensuremath{\mathcal{P}}}}1
    {∗}{{\ensuremath{\ast}}}1
    {•}{{\ensuremath{\bullet}}}1,
  % Comments
  morecomment=[l]{--},
  morecomment=[s]{/-}{-/},
  morestring=[b]",
  stringstyle=\color{stringcolor},
  keepspaces=true,
  keywordstyle=[1]\color{keywordcolor}\bfseries,
  keywordstyle=[2]\color{sortcolor}\bfseries,
  keywordstyle=[3]\color{tacticcolor}\bfseries,
  keywordstyle=[4]\color{errorcolor}\bfseries\underline,
  escapeinside={(*@}{@*)},
  columns=fullflexible,
  basicstyle=\ttfamily\small,
  commentstyle=\color{commentcolor}\itshape,
  breaklines=true,
  showstringspaces=false,
  frame=single,
  framesep=2mm,
  xleftmargin=4mm,
  numbers=left,
  numberstyle=\tiny\color{commentcolor},
  numbersep=6pt,
}

% Inline Lean code in prose: \lean{Int.ModEq.pow}
\newcommand{\lean}[1]{\lstinline[language=lean,basicstyle=\ttfamily\normalsize]|#1|}

\lstset{language=lean}

% --- local rendering patch (2026-09-30; lineage papers/cyl3-mod23 + ladder
% precedent) --- tectonic/XeTeX does not apply the listings literate table to
% multi-byte characters, so the Unicode set used by the listings of THIS note
% is mapped to math glyphs as active characters via the standard
% \lccode/\lowercase idiom (ASCII-only source, no extra packages). Characters
% deliberately NOT mapped: U+00B7 (middle dot) and U+2014 (em-dash) --- both
% occur in this note's listings and xeCJK sets them up at load time (making
% them active first aborts the run with "Control sequence ... already
% defined"; observed 2026-09-30 on this note and 2026-09-27 in the pendant
% note); they route through xeCJK + SimSun, as do all CJK ideographs and CJK
% punctuation in the Lean source comments. Rendering only: listing source
% bytes are untouched, so the byte-exactness checks are unaffected, and the
% \ifdefined guard leaves the pdflatex path identical to the template.
\ifdefined\XeTeXversion
  {\lccode`\~="2115 \lowercase{\global\catcode`~=\active \global\def~{\ensuremath{\mathbb{N}}}}}
  {\lccode`\~="2124 \lowercase{\global\catcode`~=\active \global\def~{\ensuremath{\mathbb{Z}}}}}
  {\lccode`\~="03B1 \lowercase{\global\catcode`~=\active \global\def~{\ensuremath{\alpha}}}}
  {\lccode`\~="03B2 \lowercase{\global\catcode`~=\active \global\def~{\ensuremath{\beta}}}}
  {\lccode`\~="03B9 \lowercase{\global\catcode`~=\active \global\def~{\ensuremath{\iota}}}}
  {\lccode`\~="03C6 \lowercase{\global\catcode`~=\active \global\def~{\ensuremath{\phi}}}}
  {\lccode`\~="03C8 \lowercase{\global\catcode`~=\active \global\def~{\ensuremath{\psi}}}}
  {\lccode`\~="2190 \lowercase{\global\catcode`~=\active \global\def~{\ensuremath{\leftarrow}}}}
  {\lccode`\~="2191 \lowercase{\global\catcode`~=\active \global\def~{\ensuremath{\uparrow}}}}
  {\lccode`\~="2192 \lowercase{\global\catcode`~=\active \global\def~{\ensuremath{\rightarrow}}}}
  {\lccode`\~="2194 \lowercase{\global\catcode`~=\active \global\def~{\ensuremath{\leftrightarrow}}}}
  {\lccode`\~="21A6 \lowercase{\global\catcode`~=\active \global\def~{\ensuremath{\mapsto}}}}
  {\lccode`\~="21D2 \lowercase{\global\catcode`~=\active \global\def~{\ensuremath{\Rightarrow}}}}
  {\lccode`\~="2200 \lowercase{\global\catcode`~=\active \global\def~{\ensuremath{\forall}}}}
  {\lccode`\~="2203 \lowercase{\global\catcode`~=\active \global\def~{\ensuremath{\exists}}}}
  {\lccode`\~="2205 \lowercase{\global\catcode`~=\active \global\def~{\ensuremath{\varnothing}}}}
  {\lccode`\~="2208 \lowercase{\global\catcode`~=\active \global\def~{\ensuremath{\in}}}}
  {\lccode`\~="2209 \lowercase{\global\catcode`~=\active \global\def~{\ensuremath{\notin}}}}
  {\lccode`\~="2211 \lowercase{\global\catcode`~=\active \global\def~{\ensuremath{\sum}}}}
  {\lccode`\~="2212 \lowercase{\global\catcode`~=\active \global\def~{\ensuremath{-}}}}
  {\lccode`\~="2227 \lowercase{\global\catcode`~=\active \global\def~{\ensuremath{\wedge}}}}
  {\lccode`\~="2228 \lowercase{\global\catcode`~=\active \global\def~{\ensuremath{\vee}}}}
  {\lccode`\~="2229 \lowercase{\global\catcode`~=\active \global\def~{\ensuremath{\cap}}}}
  {\lccode`\~="222A \lowercase{\global\catcode`~=\active \global\def~{\ensuremath{\cup}}}}
  {\lccode`\~="2260 \lowercase{\global\catcode`~=\active \global\def~{\ensuremath{\neq}}}}
  {\lccode`\~="2261 \lowercase{\global\catcode`~=\active \global\def~{\ensuremath{\equiv}}}}
  {\lccode`\~="2264 \lowercase{\global\catcode`~=\active \global\def~{\ensuremath{\leq}}}}
  {\lccode`\~="2265 \lowercase{\global\catcode`~=\active \global\def~{\ensuremath{\geq}}}}
  {\lccode`\~="2286 \lowercase{\global\catcode`~=\active \global\def~{\ensuremath{\subseteq}}}}
  {\lccode`\~="22A2 \lowercase{\global\catcode`~=\active \global\def~{\ensuremath{\vdash}}}}
  {\lccode`\~="00AC \lowercase{\global\catcode`~=\active \global\def~{\ensuremath{\neg}}}}
  {\lccode`\~="00B2 \lowercase{\global\catcode`~=\active \global\def~{\ensuremath{^2}}}}
  {\lccode`\~="00B3 \lowercase{\global\catcode`~=\active \global\def~{\ensuremath{^3}}}}
  {\lccode`\~="00A7 \lowercase{\global\catcode`~=\active \global\def~{\S}}}
  {\lccode`\~="00D7 \lowercase{\global\catcode`~=\active \global\def~{\ensuremath{\times}}}}
  {\lccode`\~="1D65 \lowercase{\global\catcode`~=\active \global\def~{\ensuremath{_{\mathrm{v}}}}}}
  {\lccode`\~="2081 \lowercase{\global\catcode`~=\active \global\def~{\ensuremath{_{1}}}}}
  {\lccode`\~="2082 \lowercase{\global\catcode`~=\active \global\def~{\ensuremath{_{2}}}}}
  {\lccode`\~="25B8 \lowercase{\global\catcode`~=\active \global\def~{\ensuremath{\blacktriangleright}}}}
  {\lccode`\~="25A1 \lowercase{\global\catcode`~=\active \global\def~{\ensuremath{\square}}}}
  {\lccode`\~="27E8 \lowercase{\global\catcode`~=\active \global\def~{\ensuremath{\langle}}}}
  {\lccode`\~="27E9 \lowercase{\global\catcode`~=\active \global\def~{\ensuremath{\rangle}}}}
  {\lccode`\~="27F9 \lowercase{\global\catcode`~=\active \global\def~{\ensuremath{\Longrightarrow}}}}
  {\lccode`\~="2460 \lowercase{\global\catcode`~=\active \global\def~{(1)}}}
  {\lccode`\~="2461 \lowercase{\global\catcode`~=\active \global\def~{(2)}}}
  {\lccode`\~="2462 \lowercase{\global\catcode`~=\active \global\def~{(3)}}}
\fi
\ifdefined\XeTeXversion
  \usepackage{xeCJK}
  \setCJKmainfont{SimSun}
\fi

\newtheorem{theorem}{Theorem}

\title{Independent sets of the three-row grid strip with the middle cell of
every second column deleted: a sixth-order recurrence\\ (Lean 4
machine-checked claim of record)}
\author{Aurora0134\thanks{Contact: via the GitHub account Aurora0134; ORCID
placeholder --- to be filled in by the author before any upload. No personal
information is carried in this note.}}
\date{September 30, 2026}

\begin{document}
\maketitle

\begin{abstract}
Let $z(n)$ be the number of independent sets (empty set included) of the
three-row by $n$-column grid strip $P_3\square P_n$ with the middle-row
vertex of every odd-numbered column (0-based columns $1,3,5,\dots$) deleted.
The AI4Math pipeline formalised this counting sequence in Lean 4 and
machine-proved its sixth-order recurrence: for every $n \ge 6$,
$z(n) + 15\,z(n-4) = 12\,z(n-2) + 2\,z(n-6)$, so each term is an integral
linear combination of the terms of the same parity two and four steps back.
The proof compiles with no \texttt{sorry}, and \texttt{\#print axioms}
reports exactly \{propext, Quot.sound, Classical.choice\}. The pipeline's
exclusivity review found no occupying record for the sequence $z(n)$ itself
or for the defective-strip counting object: the even-indexed subsequence of
$z$ coincides term-by-term with the registered sequence A122011 (a $3\times3$
matrix-power example with no graph-theoretic content, a boundary that the
recurrence signature is already indexed under), while $z(n)$ itself, its
odd-indexed subsequence, and the defective-strip object are unregistered.
The claim is strictly \emph{new sequence (lower tier)}, never new
mathematics and never a new recurrence. This note is a priority claim of
record, not a full paper.
\end{abstract}

\section{Introduction}

One theorem about one integer sequence is reported here. The sequence counts
independent sets of a grid strip with a periodic column defect: in the
three-row strip every odd column loses its middle vertex, so a column is
either a full $3$-vertex path or two disconnected corner vertices. The first
values are $z(0)=1, z(1)=5, z(2)=13, z(3)=47, z(4)=141, z(5)=491,
z(6)=1499, z(7)=5197$ (computed), and the theorem proved here is the
recurrence $z(n) + 15 z(n-4) = 12 z(n-2) + 2 z(n-6)$ for all $n \ge 6$. The
proof was produced by the AI4Math pipeline, in which LLM workers generate
statements and proof scripts and the Lean 4 kernel is the sole adjudicator
at every gate: the statement is frozen by hash before proving, and a separate
audit re-computes the symmetric difference, scans the whole file for
\texttt{sorry}/\texttt{admit} including comments, and re-prints the axiom
dependencies \cite{lean4}. This note is a priority claim of record deposited
so that a dated, citable public record exists ahead of the narrative paper;
no literature review is attempted here.

\section{Statement}\label{sec:statement}

The formalisation works directly with the vertex set
$\texttt{Fin }3 \times \texttt{Fin }n$: \texttt{holes} is the set of
middle-row (row 1) vertices in odd-numbered columns, and
\texttt{indepCountHoles} counts vertex subsets that are disjoint from the
hole set and independent in the box product
$\texttt{pathGraph }3\ \square\ \texttt{pathGraph }n$; the sequence is
\texttt{z n := indepCountHoles 3 n (holes 3 n)}. The recurrence is stated
with natural-number truncated subtraction, self-consistent under the
hypothesis $6 \le n$.

% LEAN: tasks/20260930-comb12-grid-indep :: c_z_rec  grade=完全证明
\begin{theorem}[sixth-order recurrence]\label{thm:rec}
For every $n : \mathbb{N}$ with $6 \le n$,
\[
\texttt{z}\,n + 15 \cdot \texttt{z}\,(n-4) = 12 \cdot \texttt{z}\,(n-2) +
2 \cdot \texttt{z}\,(n-6).
\]
\end{theorem}

\begin{lstlisting}[caption={Frozen base definitions and the head of
Theorem~\ref{thm:rec}, verbatim lines 21--48 of
\texttt{formalized/statement.lean} (sha256 \texttt{b726f0d8}): the two
decidability instances, \texttt{indepCountHoles}, \texttt{holes},
\texttt{z}, and the theorem statement. The frozen file's proof body
and its \texttt{native\_decide}
data-check lines are excluded.},label={lst:s1}]
-- 底座三件（照抄 smoke/B-deep-smokes.lean 已实测通过的写法）——

instance (n : ℕ) : DecidableRel (pathGraph n).Adj := fun _ _ =>
  decidable_of_iff _ pathGraph_adj.symm

instance {α β : Type*} [DecidableEq α] [DecidableEq β] (G : SimpleGraph α) (H : SimpleGraph β)
    [DecidableRel G.Adj] [DecidableRel H.Adj] : DecidableRel (G □ H).Adj := fun _ _ =>
  decidable_of_iff _ boxProd_adj.symm

/-- 与挖点集 `holes` 不交的独立集数 = 挖去 `holes` 后剩余图的独立集数。 -/
def indepCountHoles (m n : ℕ) (holes : Finset (Fin m × Fin n)) : ℕ :=
  ((univ : Finset (Finset (Fin m × Fin n))).filter (fun s =>
    Disjoint s holes ∧ (pathGraph m □ pathGraph n).IsIndepSet s)).card

-- 缺陷集与计数列 z ——

/-- 隔列挖中点缺陷集：`Fin m × Fin n` 中行号 = 1（中行）且列号为奇数的点，
0 基行列，即 (1, 1), (1, 3), (1, 5), …（`Finset.univ.filter` 口径，不经 `Subgraph.induce`）。 -/
def holes (m n : ℕ) : Finset (Fin m × Fin n) :=
  (univ : Finset (Fin m × Fin n)).filter (fun p => p.1.val = 1 ∧ p.2.val % 2 = 1)

/-- 三行带每隔一列挖去中行顶点所得图的独立集数列。 -/
def z (n : ℕ) : ℕ := indepCountHoles 3 n (holes 3 n)

-- 主定理（proof 占位：本任务只产 statement）——

/-- z 的六阶递推（ℕ 截断减法形，系数 15/12/2 与方向依任务书不动）。 -/
theorem c_z_rec : ∀ n : ℕ, 6 ≤ n → z n + 15 * z (n - 4) = 12 * z (n - 2) + 2 * z (n - 6) := by
\end{lstlisting}

\begin{lstlisting}[caption={Frozen attack interface and matrix certificates,
verbatim lines 18--163 of \texttt{formalized/interface.lean} (sha256
\texttt{66c96c3e}): the frozen statement segment, the profile/transfer
definitions \texttt{P5}/\texttt{Q4}/\texttt{allowed}/\texttt{ValidProf}/%
\texttt{validCount}, the matrices \texttt{Amat}/\texttt{Bmat}/\texttt{Tmat}%
\ and the two \texttt{decide}-closed certificates
\texttt{cert\_ones}/\texttt{cert\_vec0}.},label={lst:s2}]
-- ============================================================
-- 段一：冻结 statement（逐字节，禁改）
-- ============================================================

instance (n : ℕ) : DecidableRel (pathGraph n).Adj := fun _ _ =>
  decidable_of_iff _ pathGraph_adj.symm

instance {α β : Type*} [DecidableEq α] [DecidableEq β] (G : SimpleGraph α) (H : SimpleGraph β)
    [DecidableRel G.Adj] [DecidableRel H.Adj] : DecidableRel (G □ H).Adj := fun _ _ =>
  decidable_of_iff _ boxProd_adj.symm

/-- 与挖点集 `holes` 不交的独立集数 = 挖去 `holes` 后剩余图的独立集数。 -/
def indepCountHoles (m n : ℕ) (holes : Finset (Fin m × Fin n)) : ℕ :=
  ((univ : Finset (Finset (Fin m × Fin n))).filter (fun s =>
    Disjoint s holes ∧ (pathGraph m □ pathGraph n).IsIndepSet s)).card

/-- 隔列挖中点缺陷集：`Fin m × Fin n` 中行号 = 1（中行）且列号为奇数的点，
0 基行列，即 (1, 1), (1, 3), (1, 5), …（`Finset.univ.filter` 口径，不经 `Subgraph.induce`）。 -/
def holes (m n : ℕ) : Finset (Fin m × Fin n) :=
  (univ : Finset (Fin m × Fin n)).filter (fun p => p.1.val = 1 ∧ p.2.val % 2 = 1)

/-- 三行带每隔一列挖去中行顶点所得图的独立集数列。 -/
def z (n : ℕ) : ℕ := indepCountHoles 3 n (holes 3 n)

-- ============================================================
-- 段二：攻证接口（新增；证明员只读引用，不改写）
-- ============================================================

/-- 列切片：点集 `s` 在第 `j` 列上的行集合。 -/
def colSlice (n : ℕ) (j : Fin n) (s : Finset (Fin 3 × Fin n)) : Finset (Fin 3) :=
  univ.filter (fun r => (r, j) ∈ s)

/-- 偶列（无洞）允许的剖面集：全 5 个竖直独立剖面。 -/
def Pprof : Finset (Finset (Fin 3)) := {∅, {0}, {1}, {2}, {0, 2}}

/-- 奇列（有洞，禁中行）允许的剖面集：4 个不含行 1 的竖直独立剖面。 -/
def Qprof : Finset (Finset (Fin 3)) := {∅, {0}, {2}, {0, 2}}

/-- 第 `j` 列（0 基）允许的剖面集：奇列禁中行，偶列全允许。 -/
def allowed (j : ℕ) : Finset (Finset (Fin 3)) := if j % 2 = 1 then Qprof else Pprof

/-- 剖面函数合法性：每列剖面合法，且相邻两列剖面不交。 -/
def ValidProf (n : ℕ) (f : Fin n → Finset (Fin 3)) : Prop :=
  (∀ j : Fin n, f j ∈ allowed j.val) ∧
    (∀ j j' : Fin n, (pathGraph n).Adj j j' → Disjoint (f j) (f j'))

/-- `ValidProf` 的可判定性实例（显式给出，使 `validCount` 的 `filter` 实例唯一、
定义上透明：`validCount n = (univ.filter (ValidProf n)).card` 按 `rfl` 成立）。
不显式给会走 `classical` 盲选，证明里另写的同形 filter 会拿到不同实例，两侧
命题相等但不同定义，`card_bij` 直接卡死。 -/
instance (n : ℕ) : DecidablePred (ValidProf n) := fun f => by
  unfold ValidProf
  infer_instance

/-- 合法剖面函数计数。
（`noncomputable`：上面的可判定实例内部含经典判定；公理走白名单内 `Classical.choice`。） -/
noncomputable def validCount (n : ℕ) : ℕ :=
  ((univ : Finset (Fin n → Finset (Fin 3))).filter (ValidProf n)).card

/-- 由剖面函数重建点集：第 `j` 列取 `f j` 中的行。 -/
def reconstruct (n : ℕ) (f : Fin n → Finset (Fin 3)) : Finset (Fin 3 × Fin n) :=
  univ.biUnion (fun j => (f j).image (fun r => (r, j)))

-- ============================================================
-- 段三：转移矩阵与证书（`decide` 一击闭；已严格模式实测 exit 0）
-- ============================================================

/-- 剖面枚举（偶列侧，5 个）。 -/
def P5 : Fin 5 → Finset (Fin 3)
  | 0 => ∅
  | 1 => {0}
  | 2 => {1}
  | 3 => {2}
  | 4 => {0, 2}

/-- 剖面枚举（奇列侧，4 个）。 -/
def Q4 : Fin 4 → Finset (Fin 3)
  | 0 => ∅
  | 1 => {0}
  | 2 => {2}
  | 3 => {0, 2}

/-- 偶→奇转移矩阵：`Amat i j = 1` 当且仅当剖面 `P5 i` 与 `Q4 j` 不交。 -/
def Amat : Matrix (Fin 5) (Fin 4) ℤ := fun i j => if Disjoint (P5 i) (Q4 j) then 1 else 0

/-- 奇→偶转移矩阵：`Bmat i j = 1` 当且仅当剖面 `Q4 i` 与 `P5 j` 不交。 -/
def Bmat : Matrix (Fin 4) (Fin 5) ℤ := fun i j => if Disjoint (Q4 i) (P5 j) then 1 else 0

/-- 双步转移矩阵 `T = A·B`（5×5）。 -/
def Tmat : Matrix (Fin 5) (Fin 5) ℤ := Amat * Bmat

/-- 全 1 向量（5 维）。 -/
def ones5 : Fin 5 → ℤ := fun _ => 1

/-- 全 1 向量（4 维）。 -/
def ones4 : Fin 4 → ℤ := fun _ => 1

/-- 偶侧初值向量 `vec0 = A·1`（5 维；分量 = 与 `P5 i` 不交的奇剖面个数）。 -/
def vec0 : Fin 5 → ℤ := Amat *ᵥ ones4

/-- 证书一：`p(T)·1 = 0`（点态），`p(x) = x³ − 12x² + 15x − 2`。 -/
theorem cert_ones : ∀ i : Fin 5,
    ((Tmat ^ 3) *ᵥ ones5) i
      = 12 * (((Tmat ^ 2) *ᵥ ones5) i) - 15 * ((Tmat ^ 1) *ᵥ ones5) i + 2 * ones5 i := by
  decide

/-- 证书二：`p(T)·vec0 = 0`（点态）。 -/
theorem cert_vec0 : ∀ i : Fin 5,
    ((Tmat ^ 3) *ᵥ vec0) i
      = 12 * (((Tmat ^ 2) *ᵥ vec0) i) - 15 * ((Tmat ^ 1) *ᵥ vec0) i + 2 * vec0 i := by
  decide

-- ============================================================
-- 段四：名册核验（本文件编译即证下列名字存在；供证据包引用）
-- ============================================================

#check @Finset.card_bij
#check @Finset.card_biUnion
#check @Finset.card_sigma
#check @Finset.card_filter
#check @Finset.card_image_of_injOn
#check @Finset.mem_biUnion
#check @Finset.mem_image
#check @Finset.disjoint_left
#check @Finset.filter_filter
#check @Finset.sum_filter
#check @Finset.sum_boole
#check @Finset.sum_biUnion
#check @Finset.card_eq_sum_ones
#check @Fintype.card_congr
#check @Fintype.card_fin
#check @SimpleGraph.isIndepSet_iff
#check @SimpleGraph.boxProd_adj
#check @SimpleGraph.pathGraph_adj
#check @Nat.strong_induction_on
#check @Fin.last

-- 攻证期实测可用的 API 名册（2026-09-30 探针核实）——
#check @Finset.card_bij -- 依赖型签名：i : (a : α) → a ∈ s → β（带成员性证明）
#check @Fintype.card_congr -- 取 Equiv
#check @Finset.card_image_of_injOn
#check @Matrix.mulVec -- 记号 `M ᵥ* v`：(M ᵥ* v) i = ∑ j, M i j * v j
#check @Matrix.vecMul -- 记号 `v ᵥ* M`：(v ᵥ* M) j = ∑ i, v i * M i j
#check @Finset.sum_image
-- 实测不存在：`Matrix.dotProduct`（勿回喂该名字）

\end{lstlisting}

The complete final file is reproduced verbatim in Section~\ref{sec:proof}
(whole file, no excerpting); it also ships in this deposit under
\texttt{proofs/}. The frozen statement file carries the theorem body as a
\texttt{sorry} placeholder (that is its role: the formalisation department
freezes statements, the proving department replaces bodies), followed by
five \texttt{native\_decide} data-check lines for the initial values
$z(0..4) = 1, 5, 13, 47, 141$; both lie outside every listing above. In the
final file no \texttt{native\_decide} occurs in code and the initial values
are reached through the transfer-matrix route instead, so the axiom output
stays on the whitelist (Section~\ref{sec:verify}).

\section{Proof}\label{sec:proof}

The proof is a three-lane assembly, all lanes machine-checked. Lane A
(bridge) turns the graph-side count into a column-profile count: a vertex
set is sliced into per-column row sets (\texttt{colSlice}), a profile
function rebuilds a vertex set (\texttt{reconstruct}), the two maps are
mutually inverse, and \texttt{Finset.card\_bij} transports the count, giving
\texttt{z n = validCount n} (\texttt{br\_z\_eq\_validCount}), where
\texttt{validCount} counts functions assigning each column a legal profile
(the five vertically independent row sets on even columns, the four
middle-row-free ones on odd columns) with adjacent columns disjoint. Lane B
(transfer) packages the two-step recurrence of these profile counts into
matrix form: filling the last column of a width-$(m{+}2)$ strip decomposes
over the previous column's profile with weight the indicator of profile
disjointness (\texttt{tr\_fillEnd\_step}, a \texttt{card\_bij}/%
\texttt{card\_biUnion} argument), and joint induction on the step count
gives $\texttt{validCount}(2k+1) = \sum_i (T^k \mathbf{1})_i$ and
$\texttt{validCount}(2k+2) = \sum_j (T^k \vec{v}_0)_j$ with
$T = \texttt{Tmat} = \texttt{Amat}\cdot\texttt{Bmat}$, $\vec{v}_0 =
\texttt{vec0} = \texttt{Amat}\cdot\mathbf{1}_4$ (\texttt{tr\_validCount%
\_odd}, \texttt{tr\_validCount\_even}, \texttt{tr\_validCount\_zero}). Lane
C (algebra) uses the certificate polynomial $p(x) = x^3 - 12x^2 + 15x - 2$:
the two \texttt{decide}-closed certificates state $p(T)\mathbf{1} = 0$ and
$p(T)\vec{v}_0 = 0$ pointwise, which by associativity of matrix powers
yields the order-3 pointwise recurrences of $w(k) = T^k\mathbf{1}$ and
$u(k) = T^k\vec{v}_0$, hence of their sums (\texttt{asm\_w\_rec\_sum},
\texttt{asm\_u\_rec\_sum}); the main theorem splits on the parity of $n$
(the odd lane needs no boundary case; the even lane has exactly one, $n=6$,
closed by a \texttt{decide}-ed integer identity of four matrix sums,
$1499 + 15 \cdot 13 = 12 \cdot 141 + 2 \cdot 1$).

\begin{lstlisting}[basicstyle=\ttfamily\footnotesize,caption={Complete proof
file of Theorem~\ref{thm:rec}, verbatim lines 1--1188 (whole file) of
\texttt{attempts/final-c\_z\_rec.lean} (sha256 \texttt{9a837835}). Mixed
line endings on the working disk (lines 1--485 and 953--1189 LF, 486--952
CRLF) are spliced LF-normalised per the chorded/ladder precedent;
byte-exact original in \texttt{proofs/}.},label={lst:proof}]
import Mathlib

open Finset SimpleGraph Matrix

/-!
# comb-12 接口冻结件（formalized/interface.lean，2026-09-30）

任务：tasks/20260930-comb12-grid-indep（pool-comb-12 全流程）
冻结 statement 来源：tasks/20260929-preselect-deeprecheck-01/attempts/L13C/statement.lean
  （sha256 b726f0d8e01672b39a93ce0c256f5036014bc80f7f0b6a5fcb484272d805c884）

本文件 = ①冻结 statement 的声明段（逐字节）+ ②攻证接口（新增 def，kernel 实测可编）
+ ③转移矩阵证书（`decide` 一击闭，严格模式实测 exit 0）。

**接口纪律**：本文件在最终稿中逐字节保留；证明员只读引用，不得改写其中任何声明。
-/

-- ============================================================
-- 段一：冻结 statement（逐字节，禁改）
-- ============================================================

instance (n : ℕ) : DecidableRel (pathGraph n).Adj := fun _ _ =>
  decidable_of_iff _ pathGraph_adj.symm

instance {α β : Type*} [DecidableEq α] [DecidableEq β] (G : SimpleGraph α) (H : SimpleGraph β)
    [DecidableRel G.Adj] [DecidableRel H.Adj] : DecidableRel (G □ H).Adj := fun _ _ =>
  decidable_of_iff _ boxProd_adj.symm

/-- 与挖点集 `holes` 不交的独立集数 = 挖去 `holes` 后剩余图的独立集数。 -/
def indepCountHoles (m n : ℕ) (holes : Finset (Fin m × Fin n)) : ℕ :=
  ((univ : Finset (Finset (Fin m × Fin n))).filter (fun s =>
    Disjoint s holes ∧ (pathGraph m □ pathGraph n).IsIndepSet s)).card

/-- 隔列挖中点缺陷集：`Fin m × Fin n` 中行号 = 1（中行）且列号为奇数的点，
0 基行列，即 (1, 1), (1, 3), (1, 5), …（`Finset.univ.filter` 口径，不经 `Subgraph.induce`）。 -/
def holes (m n : ℕ) : Finset (Fin m × Fin n) :=
  (univ : Finset (Fin m × Fin n)).filter (fun p => p.1.val = 1 ∧ p.2.val % 2 = 1)

/-- 三行带每隔一列挖去中行顶点所得图的独立集数列。 -/
def z (n : ℕ) : ℕ := indepCountHoles 3 n (holes 3 n)

-- ============================================================
-- 段二：攻证接口（新增；证明员只读引用，不改写）
-- ============================================================

/-- 列切片：点集 `s` 在第 `j` 列上的行集合。 -/
def colSlice (n : ℕ) (j : Fin n) (s : Finset (Fin 3 × Fin n)) : Finset (Fin 3) :=
  univ.filter (fun r => (r, j) ∈ s)

/-- 偶列（无洞）允许的剖面集：全 5 个竖直独立剖面。 -/
def Pprof : Finset (Finset (Fin 3)) := {∅, {0}, {1}, {2}, {0, 2}}

/-- 奇列（有洞，禁中行）允许的剖面集：4 个不含行 1 的竖直独立剖面。 -/
def Qprof : Finset (Finset (Fin 3)) := {∅, {0}, {2}, {0, 2}}

/-- 第 `j` 列（0 基）允许的剖面集：奇列禁中行，偶列全允许。 -/
def allowed (j : ℕ) : Finset (Finset (Fin 3)) := if j % 2 = 1 then Qprof else Pprof

/-- 剖面函数合法性：每列剖面合法，且相邻两列剖面不交。 -/
def ValidProf (n : ℕ) (f : Fin n → Finset (Fin 3)) : Prop :=
  (∀ j : Fin n, f j ∈ allowed j.val) ∧
    (∀ j j' : Fin n, (pathGraph n).Adj j j' → Disjoint (f j) (f j'))

/-- `ValidProf` 的可判定性实例（显式给出，使 `validCount` 的 `filter` 实例唯一、
定义上透明：`validCount n = (univ.filter (ValidProf n)).card` 按 `rfl` 成立）。
不显式给会走 `classical` 盲选，证明里另写的同形 filter 会拿到不同实例，两侧
命题相等但不同定义，`card_bij` 直接卡死。 -/
instance (n : ℕ) : DecidablePred (ValidProf n) := fun f => by
  unfold ValidProf
  infer_instance

/-- 合法剖面函数计数。
（`noncomputable`：上面的可判定实例内部含经典判定；公理走白名单内 `Classical.choice`。） -/
noncomputable def validCount (n : ℕ) : ℕ :=
  ((univ : Finset (Fin n → Finset (Fin 3))).filter (ValidProf n)).card

/-- 由剖面函数重建点集：第 `j` 列取 `f j` 中的行。 -/
def reconstruct (n : ℕ) (f : Fin n → Finset (Fin 3)) : Finset (Fin 3 × Fin n) :=
  univ.biUnion (fun j => (f j).image (fun r => (r, j)))

-- ============================================================
-- 段三：转移矩阵与证书（`decide` 一击闭；已严格模式实测 exit 0）
-- ============================================================

/-- 剖面枚举（偶列侧，5 个）。 -/
def P5 : Fin 5 → Finset (Fin 3)
  | 0 => ∅
  | 1 => {0}
  | 2 => {1}
  | 3 => {2}
  | 4 => {0, 2}

/-- 剖面枚举（奇列侧，4 个）。 -/
def Q4 : Fin 4 → Finset (Fin 3)
  | 0 => ∅
  | 1 => {0}
  | 2 => {2}
  | 3 => {0, 2}

/-- 偶→奇转移矩阵：`Amat i j = 1` 当且仅当剖面 `P5 i` 与 `Q4 j` 不交。 -/
def Amat : Matrix (Fin 5) (Fin 4) ℤ := fun i j => if Disjoint (P5 i) (Q4 j) then 1 else 0

/-- 奇→偶转移矩阵：`Bmat i j = 1` 当且仅当剖面 `Q4 i` 与 `P5 j` 不交。 -/
def Bmat : Matrix (Fin 4) (Fin 5) ℤ := fun i j => if Disjoint (Q4 i) (P5 j) then 1 else 0

/-- 双步转移矩阵 `T = A·B`（5×5）。 -/
def Tmat : Matrix (Fin 5) (Fin 5) ℤ := Amat * Bmat

/-- 全 1 向量（5 维）。 -/
def ones5 : Fin 5 → ℤ := fun _ => 1

/-- 全 1 向量（4 维）。 -/
def ones4 : Fin 4 → ℤ := fun _ => 1

/-- 偶侧初值向量 `vec0 = A·1`（5 维；分量 = 与 `P5 i` 不交的奇剖面个数）。 -/
def vec0 : Fin 5 → ℤ := Amat *ᵥ ones4

/-- 证书一：`p(T)·1 = 0`（点态），`p(x) = x³ − 12x² + 15x − 2`。 -/
theorem cert_ones : ∀ i : Fin 5,
    ((Tmat ^ 3) *ᵥ ones5) i
      = 12 * (((Tmat ^ 2) *ᵥ ones5) i) - 15 * ((Tmat ^ 1) *ᵥ ones5) i + 2 * ones5 i := by
  decide

/-- 证书二：`p(T)·vec0 = 0`（点态）。 -/
theorem cert_vec0 : ∀ i : Fin 5,
    ((Tmat ^ 3) *ᵥ vec0) i
      = 12 * (((Tmat ^ 2) *ᵥ vec0) i) - 15 * ((Tmat ^ 1) *ᵥ vec0) i + 2 * vec0 i := by
  decide

-- ============================================================
-- 段四：名册核验（本文件编译即证下列名字存在；供证据包引用）
-- ============================================================

#check @Finset.card_bij
#check @Finset.card_biUnion
#check @Finset.card_sigma
#check @Finset.card_filter
#check @Finset.card_image_of_injOn
#check @Finset.mem_biUnion
#check @Finset.mem_image
#check @Finset.disjoint_left
#check @Finset.filter_filter
#check @Finset.sum_filter
#check @Finset.sum_boole
#check @Finset.sum_biUnion
#check @Finset.card_eq_sum_ones
#check @Fintype.card_congr
#check @Fintype.card_fin
#check @SimpleGraph.isIndepSet_iff
#check @SimpleGraph.boxProd_adj
#check @SimpleGraph.pathGraph_adj
#check @Nat.strong_induction_on
#check @Fin.last

-- 攻证期实测可用的 API 名册（2026-09-30 探针核实）——
#check @Finset.card_bij -- 依赖型签名：i : (a : α) → a ∈ s → β（带成员性证明）
#check @Fintype.card_congr -- 取 Equiv
#check @Finset.card_image_of_injOn
#check @Matrix.mulVec -- 记号 `M ᵥ* v`：(M ᵥ* v) i = ∑ j, M i j * v j
#check @Matrix.vecMul -- 记号 `v ᵥ* M`：(v ᵥ* M) j = ∑ i, v i * M i j
#check @Finset.sum_image
-- 实测不存在：`Matrix.dotProduct`（勿回喂该名字）

/-!
## 攻证须知（2026-09-30 实测，写 Lean 前必读）

1. **`decide` 无法从图定义直接算 z**：`example : z 0 = 1 := by decide` 实测失败
   （reduction 在 `instDecidableEqNat (z 0) 1` 处卡死）。z 的一切具体取值都必须走
   `z = validCount = 矩阵表达式` 路线后用 `decide` 闭（5×5 整数矩阵规模，
   `cert_ones`/`cert_vec0` 已实测 `decide` 一击闭）。
   **禁止用 `native_decide`**：它引入 `Lean.ofReduceBool`，不在公理白名单
   {propext, Quot.sound, Classical.choice} 内，入库即拒收（宪条 3）。
2. `validCount n = (univ.filter (ValidProf n)).card` 按 `rfl` 成立（见上方显式实例）；
   证明里可直接 `unfold validCount` 或 `show`。
3. `z 0 = 1` 不要从图定义硬算：走 `z 0 = validCount 0`（桥梁引理）+
   `validCount 0 = 1`（`Fin 0 → Finset (Fin 3)` 是 subsingleton，univ 恰一元素且合法）。
-/
-- ============================================================
-- 段 A · 桥梁恒等式 `z n = validCount n`（攻证员：桥梁证明员）
-- 前缀纪律：全部新引理 `br_` 开头（guide §7.2），不含接口已有声明。
-- 数学路线：guide §3（Finset 双射：独立集 ↔ 合法剖面函数）
-- ============================================================

/-!
### A0. 两侧计数对象与定义展开

`br_S1` / `br_S2` 与 `indepCountHoles` / `validCount` 的 filter 同形，
`DecidablePred` 实例链与 interface 段一/段二完全一致（`card_bij` 实例陷阱的
防法：不另写同形 filter，只复用这两个 def）。
-/

/-- 桥侧辅助：`z n` 所数的对象 = 三行 n 列图中与 `holes 3 n` 不交的独立集。 -/
def br_S1 (n : ℕ) : Finset (Finset (Fin 3 × Fin n)) :=
  (univ : Finset (Finset (Fin 3 × Fin n))).filter (fun s =>
    Disjoint s (holes 3 n) ∧ (pathGraph 3 □ pathGraph n).IsIndepSet s)

/-- 桥侧辅助：`validCount n` 所数的对象 = 合法剖面函数。 -/
def br_S2 (n : ℕ) : Finset (Fin n → Finset (Fin 3)) :=
  (univ : Finset (Fin n → Finset (Fin 3))).filter (ValidProf n)

theorem br_z_eq_card_S1 (n : ℕ) : z n = (br_S1 n).card := rfl

theorem br_validCount_eq_card_S2 (n : ℕ) : validCount n = (br_S2 n).card := rfl

theorem br_mem_S1 (n : ℕ) (s : Finset (Fin 3 × Fin n)) (hd : Disjoint s (holes 3 n))
    (hi : (pathGraph 3 □ pathGraph n).IsIndepSet s) : s ∈ br_S1 n :=
  Finset.mem_filter.mpr ⟨Finset.mem_univ s, hd, hi⟩

theorem br_mem_S2 (n : ℕ) (f : Fin n → Finset (Fin 3)) (h : ValidProf n f) :
    f ∈ br_S2 n := Finset.mem_filter.mpr ⟨Finset.mem_univ f, h⟩

/-!
### A1. 成员性小工具（列切片 / 重建 / 缺陷集 / 邻接 / 独立集改写）

全部按 `defeq` 释义展开（`have h' : <展开形> := h`），不依赖 `rw` 对
beta-可约元的匹配，规避 guide §6 之外另发现的 `rw` 匹配坑。
-/

/-- 列切片成员性：`r ∈ colSlice n j s ↔ (r, j) ∈ s`（定义直接展开）。 -/
theorem br_mem_colSlice (n : ℕ) (j : Fin n) (s : Finset (Fin 3 × Fin n)) (r : Fin 3) :
    r ∈ colSlice n j s ↔ (r, j) ∈ s := by
  constructor
  · intro h
    have h' : r ∈ univ.filter (fun r => (r, j) ∈ s) := h
    exact (Finset.mem_filter.mp h').2
  · intro h
    exact Finset.mem_filter.mpr ⟨Finset.mem_univ r, h⟩

/-- 重建成员性：`(r, j) ∈ reconstruct n f ↔ r ∈ f j`（定义直接展开）。 -/
theorem br_mem_reconstruct (n : ℕ) (f : Fin n → Finset (Fin 3)) (r : Fin 3) (j : Fin n) :
    (r, j) ∈ reconstruct n f ↔ r ∈ f j := by
  constructor
  · intro h
    have h' : (r, j) ∈ univ.biUnion (fun j' => (f j').image (fun r' => (r', j'))) := h
    obtain ⟨j', hj', himg⟩ := Finset.mem_biUnion.mp h'
    have himg' : (r, j) ∈ (f j').image (fun r' => (r', j')) := himg
    obtain ⟨r', hr', heq⟩ := Finset.mem_image.mp himg'
    obtain ⟨h1, h2⟩ := Prod.mk_inj.mp heq
    subst h2
    exact h1 ▸ hr'
  · intro hr
    have himg : (r, j) ∈ (f j).image (fun r' => (r', j)) :=
      Finset.mem_image.mpr ⟨r, hr, rfl⟩
    exact Finset.mem_biUnion.mpr ⟨j, Finset.mem_univ j, himg⟩

/-- 缺陷集成员性（中行点）。 -/
theorem br_mem_holes_iff (n : ℕ) (j : Fin n) :
    ((1 : Fin 3), j) ∈ holes 3 n ↔ j.val % 2 = 1 := by
  simp [holes, Finset.mem_filter]

/-- 缺陷集成员性（一般行）。 -/
theorem br_mem_holes_iff_row (n : ℕ) (r : Fin 3) (j : Fin n) :
    (r, j) ∈ holes 3 n ↔ r.val = 1 ∧ j.val % 2 = 1 := by
  simp [holes, Finset.mem_filter]

/-- 盒积邻接·竖直分量（逐字即 `SimpleGraph.boxProd_adj_left`）。 -/
theorem br_adj_vert {n : ℕ} {r r' : Fin 3} {j : Fin n} :
    (pathGraph 3 □ pathGraph n).Adj (r, j) (r', j) ↔ (pathGraph 3).Adj r r' :=
  boxProd_adj_left

/-- 盒积邻接·水平分量（逐字即 `SimpleGraph.boxProd_adj_right`）。 -/
theorem br_adj_horiz {n : ℕ} {r : Fin 3} {j j' : Fin n} :
    (pathGraph 3 □ pathGraph n).Adj (r, j) (r, j') ↔ (pathGraph n).Adj j j' :=
  boxProd_adj_right

/-- `IsIndepSet`（`Set.Pairwise` 形）改写为干净的 ∀∀ 否定形。
mathlib v4.34 中 `Set.Pairwise s r` 带 `x ≠ y` 前提（`Logic/Pairwise.lean:83`），
而 `SimpleGraph.ne_of_adj` 总能把该前提免费补上。 -/
theorem br_isIndepSet_iff_forall {m n : ℕ} (s : Finset (Fin m × Fin n)) :
    (pathGraph m □ pathGraph n).IsIndepSet s ↔
      ∀ a b : Fin m × Fin n, a ∈ s → b ∈ s → ¬(pathGraph m □ pathGraph n).Adj a b := by
  constructor
  · intro h a b ha hb habj
    exact h (Finset.mem_coe.mpr ha) (Finset.mem_coe.mpr hb) (ne_of_adj _ habj) habj
  · intro h
    exact fun a ha b hb _ habj => h a b (Finset.mem_coe.mp ha) (Finset.mem_coe.mp hb) habj

/-!
### A2. 竖直独立剖面的有限枚举（`Pprof` / `Qprof` 小引理）

`Finset (Fin 3)` 只有 8 个子集；枚举入口用 `fin_cases`（实测 `interval_cases`
不支持 `Fin 3`，见 notes）与 5/4 路 `Finset.mem_insert` 链。
-/

/-- 不含中行 ⇒ 属于奇列允许集 `Qprof`。 -/
theorem br_Qprof_of_not_mem_one (t : Finset (Fin 3)) (h : (1 : Fin 3) ∉ t) :
    t ∈ Qprof := by
  have hsub : t ⊆ ({0, 2} : Finset (Fin 3)) := by
    intro r hr
    fin_cases r
    · exact Finset.mem_insert.2 (Or.inl rfl)
    · exact absurd hr h
    · exact Finset.mem_insert.2 (Or.inr (Finset.mem_singleton.2 rfl))
  have hps : ({0, 2} : Finset (Fin 3)).powerset = Qprof := by decide
  rw [← hps]
  exact Finset.mem_powerset.mpr hsub

/-- 奇列允许集的成员不含中行。 -/
theorem br_Qprof_not_mem_one (t : Finset (Fin 3)) (h : t ∈ Qprof) : (1 : Fin 3) ∉ t := by
  rw [Qprof, Finset.mem_insert, Finset.mem_insert, Finset.mem_insert,
    Finset.mem_singleton] at h
  rcases h with h | h | h | h <;> subst h <;> decide

/-- 列内无竖直相邻 ⇒ 属于偶列允许集 `Pprof`。 -/
theorem br_Pprof_of_indep (t : Finset (Fin 3))
    (h : ∀ r r' : Fin 3, (pathGraph 3).Adj r r' → r ∈ t → r' ∈ t → False) :
    t ∈ Pprof := by
  by_cases h1 : (1 : Fin 3) ∈ t
  · have h0 : (0 : Fin 3) ∉ t := fun h0' => h 0 1 (by decide) h0' h1
    have h2 : (2 : Fin 3) ∉ t := fun h2' => h 1 2 (by decide) h1 h2'
    have hsub : t ⊆ ({(1 : Fin 3)} : Finset (Fin 3)) := by
      intro r hr
      fin_cases r
      · exact absurd hr h0
      · exact Finset.mem_singleton.2 rfl
      · exact absurd hr h2
    have hsub' : ({(1 : Fin 3)} : Finset (Fin 3)) ⊆ t := by
      intro r hr
      have hr1 : r = (1 : Fin 3) := Finset.mem_singleton.1 hr
      rw [hr1]
      exact h1
    have hcard : t = {(1 : Fin 3)} := Finset.Subset.antisymm hsub hsub'
    rw [hcard]
    decide
  · have hQ : t ∈ Qprof := br_Qprof_of_not_mem_one t h1
    have hQP : Qprof ⊆ Pprof := by decide
    exact hQP hQ

/-- 允许剖面（`Pprof` 成员）竖直无相邻：`A14` 偶列支的唯一入口。 -/
theorem br_Pprof_vert_indep (t : Finset (Fin 3)) (h : t ∈ Pprof) (r r' : Fin 3)
    (hadj : (pathGraph 3).Adj r r') (hr : r ∈ t) (hr' : r' ∈ t) : False := by
  have hcases : t = ∅ ∨ t = {(0 : Fin 3)} ∨ t = {(1 : Fin 3)} ∨ t = {(2 : Fin 3)}
      ∨ t = {(0 : Fin 3), (2 : Fin 3)} := by
    rw [Pprof, Finset.mem_insert, Finset.mem_insert, Finset.mem_insert, Finset.mem_insert,
      Finset.mem_singleton] at h
    tauto
  rcases hcases with h | h | h | h | h <;> subst h
  · simp at hr
  · rw [Finset.mem_singleton] at hr hr'
    subst hr
    subst hr'
    exact absurd hadj (by decide)
  · rw [Finset.mem_singleton] at hr hr'
    subst hr
    subst hr'
    exact absurd hadj (by decide)
  · rw [Finset.mem_singleton] at hr hr'
    subst hr
    subst hr'
    exact absurd hadj (by decide)
  · rw [Finset.mem_insert, Finset.mem_singleton] at hr hr'
    rcases hr with hr | hr <;> rcases hr' with hr' | hr' <;>
      subst hr <;> subst hr' <;> exact absurd hadj (by decide)

/-!
### A3. 奇偶列的允许集
-/

theorem br_allowed_odd (j : ℕ) (h : j % 2 = 1) : allowed j = Qprof := by
  simp [allowed, h]

theorem br_allowed_even (j : ℕ) (h : j % 2 ≠ 1) : allowed j = Pprof := by
  simp [allowed, h]

/-!
### A4. 双射两条互逆方向（guide §3.1 条 1、2）
-/

/-- 桥梁方向一：`ψ (φ s) = s`，无需任何假设（guide §3.1 条 1）。 -/
theorem br_reconstruct_colSlice (n : ℕ) (s : Finset (Fin 3 × Fin n)) :
    reconstruct n (fun j => colSlice n j s) = s := by
  apply Finset.ext
  intro p
  obtain ⟨r, j⟩ := p
  constructor
  · intro h
    exact (br_mem_colSlice n j s r).mp
      ((br_mem_reconstruct n (fun j => colSlice n j s) r j).mp h)
  · intro h
    exact (br_mem_reconstruct n (fun j => colSlice n j s) r j).mpr
      ((br_mem_colSlice n j s r).mpr h)

/-- 桥梁方向二：`φ (ψ f) = f`，只用到 `Prod.mk.injEq`，不需要 loopless（guide §3.1 条 2）。 -/
theorem br_colSlice_reconstruct (n : ℕ) (f : Fin n → Finset (Fin 3)) (j : Fin n) :
    colSlice n j (reconstruct n f) = f j := by
  apply Finset.ext
  intro r
  exact (br_mem_colSlice n j (reconstruct n f) r).trans (br_mem_reconstruct n f r j)

/-!
### A5. 桥梁方向三、四（合法剖面 ↔ 独立集）
-/

/-- 桥梁方向三：与 holes 不交的独立集 ⇒ 合法剖面（guide §3.1 条 3）。 -/
theorem br_ValidProf_of_mem_S1 (n : ℕ) (s : Finset (Fin 3 × Fin n)) (h : s ∈ br_S1 n) :
    ValidProf n (fun j => colSlice n j s) := by
  obtain ⟨hd, hi⟩ := (Finset.mem_filter.mp h).2
  have hi' : ∀ a b : Fin 3 × Fin n, a ∈ s → b ∈ s →
      ¬(pathGraph 3 □ pathGraph n).Adj a b := (br_isIndepSet_iff_forall s).mp hi
  refine ⟨?_, ?_⟩
  · intro j
    show colSlice n j s ∈ allowed j.val
    by_cases hodd : j.val % 2 = 1
    · rw [br_allowed_odd j.val hodd]
      refine br_Qprof_of_not_mem_one _ ?_
      intro h1
      have hmem : ((1 : Fin 3), j) ∈ s := (br_mem_colSlice n j s 1).mp h1
      have hnot : ((1 : Fin 3), j) ∉ holes 3 n := Finset.disjoint_left.mp hd hmem
      exact hnot ((br_mem_holes_iff n j).mpr hodd)
    · rw [br_allowed_even j.val hodd]
      refine br_Pprof_of_indep _ ?_
      intro r r' hadj hr hr'
      exact hi' (r, j) (r', j) ((br_mem_colSlice n j s r).mp hr)
        ((br_mem_colSlice n j s r').mp hr') ((br_adj_vert).mpr hadj)
  · intro j j' hadj
    rw [Finset.disjoint_left]
    intro r hrj hrj'
    exact hi' (r, j) (r, j') ((br_mem_colSlice n j s r).mp hrj)
      ((br_mem_colSlice n j' s r).mp hrj') ((br_adj_horiz).mpr hadj)

/-- 桥梁方向四：合法剖面 ⇒ 与 holes 不交的独立集（guide §3.1 条 4）。 -/
theorem br_mem_S1_of_ValidProf (n : ℕ) (f : Fin n → Finset (Fin 3)) (h : ValidProf n f) :
    reconstruct n f ∈ br_S1 n := by
  obtain ⟨hcol, hadjcol⟩ := h
  refine br_mem_S1 n (reconstruct n f) ?_ ?_
  · rw [Finset.disjoint_left]
    intro p hp
    obtain ⟨r, j⟩ := p
    rw [br_mem_reconstruct] at hp
    by_contra hcon
    rw [br_mem_holes_iff_row] at hcon
    obtain ⟨hval, hodd⟩ := hcon
    have hr1 : r = (1 : Fin 3) := by
      have : r.val = ((1 : Fin 3)).val := by rw [hval]; rfl
      exact Fin.val_injective this
    subst hr1
    have hcolj : f j ∈ allowed j.val := hcol j
    rw [br_allowed_odd j.val hodd] at hcolj
    exact br_Qprof_not_mem_one _ hcolj hp
  · rw [br_isIndepSet_iff_forall]
    intro a b ha hb habj
    obtain ⟨ra, ja⟩ := a
    obtain ⟨rb, jb⟩ := b
    rw [br_mem_reconstruct] at ha hb
    rcases boxProd_adj.mp habj with ⟨hvert, hj⟩ | ⟨hhoriz, hr⟩
    · -- 竖直邻接：同列相邻行
      have hvert' : (pathGraph 3).Adj ra rb := hvert
      have hj' : ja = jb := hj
      rw [← hj'] at hb
      have hP : f ja ∈ Pprof := by
        have hcolj : f ja ∈ allowed ja.val := hcol ja
        by_cases hodd : ja.val % 2 = 1
        · rw [br_allowed_odd ja.val hodd] at hcolj
          have hQP : Qprof ⊆ Pprof := by decide
          exact hQP hcolj
        · rw [br_allowed_even ja.val hodd] at hcolj
          exact hcolj
      exact br_Pprof_vert_indep _ hP ra rb hvert' ha hb
    · -- 水平邻接：同行相邻列
      have hhoriz' : (pathGraph n).Adj ja jb := hhoriz
      have hr' : ra = rb := hr
      rw [hr'] at ha
      have hdis : Disjoint (f ja) (f jb) := hadjcol ja jb hhoriz'
      exact Finset.disjoint_left.mp hdis ha hb

/-!
### A6. 计数相等（`Finset.card_bij`，guide §3.2）
-/

/-- 桥梁主引理：`z n = validCount n`（独立集 ↔ 合法剖面函数的双射）。 -/
theorem br_z_eq_validCount (n : ℕ) : z n = validCount n := by
  rw [br_z_eq_card_S1, br_validCount_eq_card_S2]
  refine Finset.card_bij (fun s _ => fun j => colSlice n j s) ?_ ?_ ?_
  · intro s hs
    exact br_mem_S2 n _ (br_ValidProf_of_mem_S1 n s hs)
  · intro s1 hs1 s2 hs2 heq
    have h1 : reconstruct n (fun j => colSlice n j s1)
        = reconstruct n (fun j => colSlice n j s2) := congrArg (reconstruct n) heq
    rw [br_reconstruct_colSlice, br_reconstruct_colSlice] at h1
    exact h1
  · intro b hb
    have hbv : ValidProf n b := (Finset.mem_filter.mp hb).2
    refine ⟨reconstruct n b, br_mem_S1_of_ValidProf n b hbv, ?_⟩
    funext j
    exact br_colSlice_reconstruct n b j

-- ============================================================
-- 段 B：转移 validCount n = 矩阵表达式（guide §4）
-- ============================================================

/-- 转移侧辅助：前 `m + 1` 列合法填法数，且最后一列剖面恰为 `p`（guide §4.2）。 -/
def tr_fillEnd : ℕ → Finset (Fin 3) → ℕ
  | 0, _ => 0
  | m + 1, p =>
    ((univ : Finset (Fin (m + 1) → Finset (Fin 3))).filter
      (fun f => ValidProf (m + 1) f ∧ f (Fin.last m) = p)).card

/-- 转移侧辅助：偶侧向量 `Vv k = Tmat ^ k *ᵥ ones5` 的递归形式（guide §4.2；
记号按 §6 用 `*ᵥ`（mulVec，矩阵在左），guide 正文的 `ᵥ*` 为同义简写）。 -/
def tr_Vv : ℕ → Fin 5 → ℤ
  | 0 => fun _ => 1
  | k + 1 => Tmat *ᵥ tr_Vv k

/-! ### Fin 列指标的分解小引理 -/

/-- `Fin (n + 1)` 的元素要么是某个 `Fin n` 元素的 `castSucc`，要么是 `Fin.last n`。 -/
theorem tr_fin_cases {n : ℕ} (i : Fin (n + 1)) :
    (∃ j : Fin n, i = j.castSucc) ∨ i = Fin.last n := by
  rcases Nat.lt_or_ge i.val n with h | h
  · left
    exact ⟨⟨i.val, h⟩, Fin.ext (by simp)⟩
  · right
    exact Fin.ext (by simp only [Fin.val_last]; have := i.isLt; omega)

/-- `pathGraph` 邻接在 `castSucc` 下保持。 -/
theorem tr_adj_castSucc {n : ℕ} (i j : Fin n) :
    (pathGraph (n + 1)).Adj (i.castSucc) (j.castSucc) ↔ (pathGraph n).Adj i j := by
  rw [pathGraph_adj, pathGraph_adj, Fin.val_castSucc, Fin.val_castSucc]

/-- `Fin (m + 2)` 的末两列相邻。 -/
theorem tr_adj_last_pair {m : ℕ} :
    (pathGraph (m + 1 + 1)).Adj ((Fin.last m).castSucc) (Fin.last (m + 1)) := by
  rw [pathGraph_adj, Fin.val_castSucc]
  simp only [Fin.val_last]
  exact Or.inl trivial

/-- `Fin (m + 2)` 中与末列相邻的只能是倒数第二列。 -/
theorem tr_adj_last_neighbor {m : ℕ} {i : Fin (m + 1)}
    (h : (pathGraph (m + 1 + 1)).Adj (i.castSucc) (Fin.last (m + 1))) : i = Fin.last m := by
  rw [pathGraph_adj, Fin.val_castSucc] at h
  simp only [Fin.val_last] at h
  have hi := i.isLt
  rcases h with h1 | h1
  · exact Fin.ext (by simp only [Fin.val_last]; omega)
  · exfalso; omega

/-- `Fin 1` 的唯一元素就是 `Fin.last 0`。 -/
theorem tr_fin_one_eq (i : Fin (0 + 1)) : i = Fin.last 0 :=
  Fin.ext (by simp only [Fin.val_last]; have := i.isLt; omega)

/-! ### B1 / B2：剖面枚举 -/

/-- B1a：`P5` 枚举恰为 Pprof。 -/
theorem tr_P5_image : Finset.image P5 univ = Pprof := by
  decide

/-- B1b：`P5` 单射。 -/
theorem tr_P5_inj : Function.Injective P5 := by
  decide

/-- B2a：`Q4` 枚举恰为 Qprof。 -/
theorem tr_Q4_image : Finset.image Q4 univ = Qprof := by
  decide

/-- B2b：`Q4` 单射。 -/
theorem tr_Q4_inj : Function.Injective Q4 := by
  decide

/-! ### B3：allowed 的奇偶 -/

/-- B3a：偶列允许集 = Pprof。 -/
theorem tr_allowed_double (j : ℕ) : allowed (2 * j) = Pprof := by
  simp [allowed]

/-- B3b：奇列允许集 = Qprof。 -/
theorem tr_allowed_double_succ (j : ℕ) : allowed (2 * j + 1) = Qprof := by
  simp [allowed]

/-! ### B4：单列基例 -/

/-- `ValidProf 1` 只约束唯一那一列。 -/
theorem tr_ValidProf_one (f : Fin (0 + 1) → Finset (Fin 3)) :
    ValidProf (0 + 1) f ↔ f (Fin.last 0) ∈ allowed 0 := by
  constructor
  · intro h
    exact h.1 (Fin.last 0)
  · intro h
    refine ⟨fun j => ?_, fun j j' hj => ?_⟩
    · have hji := tr_fin_one_eq j
      rw [hji]
      exact h
    · exfalso
      rw [pathGraph_adj] at hj
      have hj1 := j.isLt
      have hj2 := j'.isLt
      omega

/-- B4：单列递推基例。 -/
theorem tr_fillEnd_one (p : Finset (Fin 3)) :
    tr_fillEnd 1 p = if p ∈ allowed 0 then 1 else 0 := by
  by_cases hp : p ∈ allowed 0
  · have hif : (if p ∈ allowed 0 then 1 else 0) = 1 := by simp [hp]
    rw [hif]
    have h1 : tr_fillEnd 1 p
        = ((univ : Finset (Fin (0 + 1) → Finset (Fin 3))).filter
            (fun f => ValidProf (0 + 1) f ∧ f (Fin.last 0) = p)).card := rfl
    rw [h1]
    have h2 : ((univ : Finset (Fin (0 + 1) → Finset (Fin 3))).filter
        (fun f => ValidProf (0 + 1) f ∧ f (Fin.last 0) = p))
        = {fun _ : Fin (0 + 1) => p} := by
      ext f
      simp only [Finset.mem_filter, Finset.mem_univ, true_and, Finset.mem_singleton]
      rw [tr_ValidProf_one]
      constructor
      · rintro ⟨h1, h2⟩
        funext i
        have hii := tr_fin_one_eq i
        rw [hii]
        exact h2
      · rintro rfl
        exact ⟨hp, rfl⟩
    rw [h2, Finset.card_singleton]
  · have hif : (if p ∈ allowed 0 then 1 else 0) = 0 := by simp [hp]
    rw [hif]
    have h1 : tr_fillEnd 1 p
        = ((univ : Finset (Fin (0 + 1) → Finset (Fin 3))).filter
            (fun f => ValidProf (0 + 1) f ∧ f (Fin.last 0) = p)).card := rfl
    rw [h1]
    have h2 : ((univ : Finset (Fin (0 + 1) → Finset (Fin 3))).filter
        (fun f => ValidProf (0 + 1) f ∧ f (Fin.last 0) = p)) = ∅ := by
      ext f
      simp only [Finset.mem_filter, Finset.mem_univ, true_and, Finset.notMem_empty]
      exact ⟨fun h => hp (h.2 ▸ h.1.1 (Fin.last 0)), fun h => h.elim⟩
    rw [h2, Finset.card_empty]

/-! ### B5：填法数的双步递推 -/

/-- 砍掉末列：`ValidProf (m + 2) f` 给出前 `m + 1` 列的 `ValidProf (m + 1) (Fin.init f)`。 -/
theorem tr_validProf_init {m : ℕ} (f : Fin (m + 1 + 1) → Finset (Fin 3))
    (h : ValidProf (m + 1 + 1) f) : ValidProf (m + 1) (Fin.init f) := by
  refine ⟨fun i => ?_, fun i i' hi => ?_⟩
  · exact h.1 (i.castSucc)
  · exact h.2 (i.castSucc) (i'.castSucc) ((tr_adj_castSucc i i').2 hi)

/-- 接上末列：`Fin.snoc g p` 合法当且仅当 `g` 合法、`p` 在末列允许集、且与倒数第二列不交。 -/
theorem tr_validProf_snoc {m : ℕ} (g : Fin (m + 1) → Finset (Fin 3)) (p : Finset (Fin 3))
    (hg : ValidProf (m + 1) g) (hp : p ∈ allowed (m + 1))
    (hdisj : Disjoint (g (Fin.last m)) p) :
    ValidProf (m + 1 + 1) (Fin.snoc (α := fun _ => Finset (Fin 3)) g p) := by
  refine ⟨fun j => ?_, fun j j' hadj => ?_⟩
  · rcases tr_fin_cases j with ⟨i, rfl⟩ | hlast
    · rw [Fin.snoc_castSucc]
      exact hg.1 i
    · rw [hlast, Fin.snoc_last, Fin.val_last]
      exact hp
  · rcases tr_fin_cases j with ⟨i, rfl⟩ | hj
    · rcases tr_fin_cases j' with ⟨i', rfl⟩ | hj'
      · rw [Fin.snoc_castSucc, Fin.snoc_castSucc]
        exact hg.2 i i' ((tr_adj_castSucc i i').mp hadj)
      · rw [hj'] at hadj
        rw [hj', Fin.snoc_last, Fin.snoc_castSucc]
        have hi : i = Fin.last m := tr_adj_last_neighbor hadj
        rw [hi]
        exact hdisj
    · rcases tr_fin_cases j' with ⟨i', rfl⟩ | hj'
      · rw [hj] at hadj
        rw [hj, Fin.snoc_last, Fin.snoc_castSucc]
        have hi' : i' = Fin.last m := tr_adj_last_neighbor (Adj.symm hadj)
        rw [hi']
        refine Finset.disjoint_left.mpr fun r hr hr' => ?_
        exact Finset.disjoint_left.mp hdisj hr' hr
      · exfalso
        rw [hj, hj', pathGraph_adj] at hadj
        simp only [Fin.val_last] at hadj
        omega

/-- B5：填法数的双步递推（先走奇列 Amat 加权，再走偶列 Bmat 加权）。 -/
theorem tr_fillEnd_step (m : ℕ) (p : Finset (Fin 3)) :
    tr_fillEnd (m + 2) p
      = if p ∈ allowed (m + 1) then
          ∑ q ∈ allowed m, (if Disjoint q p then 1 else 0) * tr_fillEnd (m + 1) q
        else 0 := by
  have hL : tr_fillEnd (m + 2) p
      = ((univ : Finset (Fin (m + 1 + 1) → Finset (Fin 3))).filter
          (fun f => ValidProf (m + 1 + 1) f ∧ f (Fin.last (m + 1)) = p)).card := rfl
  rw [hL]
  split_ifs with hp
  · -- 主情形：末列剖面 p 合法
    have hcard : ((univ : Finset (Fin (m + 1 + 1) → Finset (Fin 3))).filter
        (fun f => ValidProf (m + 1 + 1) f ∧ f (Fin.last (m + 1)) = p)).card
        = ((allowed m).biUnion (fun q => (univ : Finset (Fin (m + 1) → Finset (Fin 3))).filter
            (fun g => ValidProf (m + 1) g ∧ g (Fin.last m) = q ∧ Disjoint q p))).card := by
      apply Finset.card_bij (i := fun f _ => Fin.init f)
      · -- 映入
        intro f hf
        obtain ⟨-, hval, hfp⟩ := Finset.mem_filter.mp hf
        refine Finset.mem_biUnion.mpr ⟨f ((Fin.last m).castSucc), ?_, ?_⟩
        · exact hval.1 ((Fin.last m).castSucc)
        · refine Finset.mem_filter.mpr ⟨Finset.mem_univ _, ?_⟩
          refine ⟨tr_validProf_init f hval, rfl, ?_⟩
          exact hfp ▸ hval.2 ((Fin.last m).castSucc) (Fin.last (m + 1)) tr_adj_last_pair
      · -- 单射
        intro f₁ hf₁ f₂ hf₂ h
        obtain ⟨-, hv1, hl1⟩ := Finset.mem_filter.mp hf₁
        obtain ⟨-, hv2, hl2⟩ := Finset.mem_filter.mp hf₂
        calc f₁ = Fin.snoc (α := fun _ => Finset (Fin 3)) (Fin.init f₁) (f₁ (Fin.last (m + 1))) :=
              (Fin.snoc_init_self f₁).symm
          _ = Fin.snoc (α := fun _ => Finset (Fin 3)) (Fin.init f₂) (f₁ (Fin.last (m + 1))) := by
              rw [h]
          _ = Fin.snoc (α := fun _ => Finset (Fin 3)) (Fin.init f₂) (f₂ (Fin.last (m + 1))) := by
              rw [hl1, hl2]
          _ = f₂ := Fin.snoc_init_self f₂
      · -- 满射
        intro b hb
        obtain ⟨q, hqm, hbf⟩ := Finset.mem_biUnion.mp hb
        obtain ⟨-, hbval, hblast, hbdisj⟩ := Finset.mem_filter.mp hbf
        refine ⟨Fin.snoc (α := fun _ => Finset (Fin 3)) b p, ?_, ?_⟩
        · refine Finset.mem_filter.mpr ⟨Finset.mem_univ _, ?_⟩
          refine ⟨tr_validProf_snoc b p hbval hp (hblast ▸ hbdisj),
            Fin.snoc_last (α := fun _ => Finset (Fin 3)) p b⟩
        · exact Fin.init_snoc (α := fun _ => Finset (Fin 3)) p b
    rw [hcard, Finset.card_biUnion]
    · refine Finset.sum_congr rfl fun q hq => ?_
      split_ifs with hd
      · rw [one_mul]
        have hfib : (univ : Finset (Fin (m + 1) → Finset (Fin 3))).filter
            (fun g => ValidProf (m + 1) g ∧ g (Fin.last m) = q ∧ Disjoint q p)
            = (univ : Finset (Fin (m + 1) → Finset (Fin 3))).filter
              (fun g => ValidProf (m + 1) g ∧ g (Fin.last m) = q) := by
          ext g
          simp only [Finset.mem_filter, Finset.mem_univ, true_and, hd, and_true]
        rw [hfib]
        rfl
      · have hfib : (univ : Finset (Fin (m + 1) → Finset (Fin 3))).filter
            (fun g => ValidProf (m + 1) g ∧ g (Fin.last m) = q ∧ Disjoint q p) = ∅ := by
          ext g
          simp only [Finset.mem_filter, Finset.mem_univ, hd, and_false]
          exact ⟨fun h => h.elim, fun h => Finset.notMem_empty g h⟩
        rw [hfib, Finset.card_empty, zero_mul]
    · -- 纤维两两不交
      intro a ha b hb hab x hxa hxb y hy
      exfalso
      apply hab
      have h1 : y (Fin.last m) = a := (((Finset.mem_filter.mp (hxa hy)).right).right).left
      have h2 : y (Fin.last m) = b := (((Finset.mem_filter.mp (hxb hy)).right).right).left
      rw [← h1, ← h2]
  · -- 末列剖面不允许：filter 为空
    have hempty : (univ : Finset (Fin (m + 1 + 1) → Finset (Fin 3))).filter
        (fun f => ValidProf (m + 1 + 1) f ∧ f (Fin.last (m + 1)) = p) = ∅ := by
      ext f
      simp only [Finset.mem_filter, Finset.mem_univ, true_and, Finset.notMem_empty, iff_false]
      rintro ⟨hval, hfp⟩
      exact hp (hfp ▸ hval.1 (Fin.last (m + 1)))
    rw [hempty, Finset.card_empty]

/-! ### B6：零列恰一合法填法 -/

/-- B6：零列恰一合法填法。 -/
theorem tr_validCount_zero : validCount 0 = 1 := by
  have hval : ∀ f : Fin 0 → Finset (Fin 3), ValidProf 0 f := fun f =>
    ⟨fun j => Fin.elim0 j, fun j j' _ => Fin.elim0 j⟩
  have hfil : (univ : Finset (Fin 0 → Finset (Fin 3))).filter (ValidProf 0) = univ := by
    ext f
    simp [hval]
  have hsub : Subsingleton (Fin 0 → Finset (Fin 3)) :=
    Subsingleton.intro fun a b => by
      funext j
      exact Fin.elim0 j
  have hcard : (univ : Finset (Fin 0 → Finset (Fin 3))).card = 1 := by
    have h1 : (univ : Finset (Fin 0 → Finset (Fin 3))).card ≤ 1 :=
      Finset.card_le_one.mpr fun a _ b _ => Subsingleton.elim a b
    have h2 : 0 < (univ : Finset (Fin 0 → Finset (Fin 3))).card :=
      Finset.card_pos.mpr Finset.univ_nonempty
    omega
  simp only [validCount]
  rw [hfil, hcard]

/-! ### B7：合法剖面数 = 按末列剖面求和 -/

/-- B7：合法剖面数 = 按末列剖面求和。 -/
theorem tr_validCount_succ (m : ℕ) :
    validCount (m + 1) = ∑ p ∈ allowed m, tr_fillEnd (m + 1) p := by
  have hfib : (univ : Finset (Fin (m + 1) → Finset (Fin 3))).filter (ValidProf (m + 1))
      = (allowed m).biUnion (fun p => (univ : Finset (Fin (m + 1) → Finset (Fin 3))).filter
          (fun f => ValidProf (m + 1) f ∧ f (Fin.last m) = p)) := by
    ext f
    simp only [Finset.mem_filter, Finset.mem_biUnion, Finset.mem_univ, true_and]
    constructor
    · intro h
      refine ⟨f (Fin.last m), h.1 (Fin.last m), h, rfl⟩
    · rintro ⟨p, -, hf, hfp⟩
      exact hf
  rw [validCount, hfib]
  rw [Finset.card_biUnion]
  · rfl
  · intro a ha b hb hab x hxa hxb y hy
    exfalso
    apply hab
    have h1 : y (Fin.last m) = a := ((Finset.mem_filter.mp (hxa hy)).right).right
    have h2 : y (Fin.last m) = b := ((Finset.mem_filter.mp (hxb hy)).right).right
    rw [← h1, ← h2]

/-! ### B8 / B9 的辅助：枚举成员性、cast 处理、if 权重对齐、联合归纳 -/

/-- 用可判定条件的证明化简 `if`（取真支）。 -/
theorem tr_if_pos {α : Type} {c : Prop} [Decidable c] (h : c) (a b : α) :
    (if c then a else b) = a := by
  simp [h]

/-- ℕ → ℤ 的 cast 穿过 `if`。 -/
theorem tr_cast_if {c : Prop} [Decidable c] (a b : ℕ) :
    ((if c then a else b : ℕ) : ℤ) = if c then (a : ℤ) else (b : ℤ) := by
  by_cases h : c
  · simp [h]
  · simp [h]

/-- ℕ → ℤ 的 cast 穿过 `(if c then 1 else 0) * f`（一步到位）。 -/
theorem tr_cast_ite_mul {c : Prop} [Decidable c] (f : ℕ) :
    (((if c then 1 else 0) * f : ℕ) : ℤ) = (if c then (1 : ℤ) else 0) * (f : ℤ) := by
  by_cases h : c
  · simp [h]
  · simp [h]

/-- `P5` 的每个分量都在 `Pprof` 中。 -/
theorem tr_P5_mem (i : Fin 5) : P5 i ∈ Pprof := by
  rw [← tr_P5_image]
  exact Finset.mem_image.mpr ⟨i, Finset.mem_univ i, rfl⟩

/-- `Q4` 的每个分量都在 `Qprof` 中。 -/
theorem tr_Q4_mem (j : Fin 4) : Q4 j ∈ Qprof := by
  rw [← tr_Q4_image]
  exact Finset.mem_image.mpr ⟨j, Finset.mem_univ j, rfl⟩

/-- `P5` 在 `univ` 上单射（`sum_image` 用）。 -/
theorem tr_P5_injOn : Set.InjOn P5 (↑(univ : Finset (Fin 5))) := by
  intro a _ b _ h
  exact tr_P5_inj h

/-- `Q4` 在 `univ` 上单射（`sum_image` 用）。 -/
theorem tr_Q4_injOn : Set.InjOn Q4 (↑(univ : Finset (Fin 4))) := by
  intro a _ b _ h
  exact tr_Q4_inj h

/-- 不交 `if` 权重交换两参数后相同（ℤ 系数版）。 -/
theorem tr_ite_disjoint_comm (a b : Finset (Fin 3)) :
    (if Disjoint a b then (1 : ℤ) else 0) = (if Disjoint b a then (1 : ℤ) else 0) := by
  by_cases h : Disjoint a b
  · have h' : Disjoint b a := _root_.disjoint_comm.mp h
    simp [h, h']
  · have h' : ¬ Disjoint b a := fun h' => h (_root_.disjoint_comm.mpr h')
    simp [h, h']

/-- `Amat i j = Bmat j i`（不交条件交换两个剖面）。 -/
theorem tr_Amat_eq_Bmat (i : Fin 5) (j : Fin 4) : Amat i j = Bmat j i := by
  show (if Disjoint (P5 i) (Q4 j) then (1 : ℤ) else 0)
    = (if Disjoint (Q4 j) (P5 i) then (1 : ℤ) else 0)
  exact tr_ite_disjoint_comm (P5 i) (Q4 j)

/-- 矩阵幂前移一步：`Tmat ^ (k + 1) = Tmat * Tmat ^ k`。 -/
theorem tr_pow_succ' (k : ℕ) : Tmat ^ (k + 1) = Tmat * Tmat ^ k := by
  rw [pow_succ]
  exact (Commute.eq ((Commute.refl Tmat).pow_right k)).symm

/-- B8/B9 的联合归纳形式（计数提升到 ℤ 后陈述）。 -/
theorem tr_fillEnd_pair (k : ℕ) :
    (∀ i : Fin 5, (tr_fillEnd (2 * k + 1) (P5 i) : ℤ) = ((Tmat ^ k) *ᵥ ones5) i) ∧
    (∀ j : Fin 4, (tr_fillEnd (2 * k + 2) (Q4 j) : ℤ) = (Bmat *ᵥ ((Tmat ^ k) *ᵥ ones5)) j) := by
  induction k with
  | zero =>
    refine ⟨?_, ?_⟩
    · intro i
      rw [show 2 * 0 + 1 = 1 from by simp, tr_fillEnd_one, tr_cast_if,
        show allowed 0 = Pprof from by simp [allowed], tr_if_pos (tr_P5_mem i)]
      rw [pow_zero, Matrix.one_mulVec, Nat.cast_one]
      rfl
    · intro j
      rw [show 2 * 0 + 2 = 0 + 2 from by simp, tr_fillEnd_step, tr_cast_if,
        show (0 : ℕ) + 1 = 1 from by simp, show allowed 1 = Qprof from by simp [allowed],
        show allowed 0 = Pprof from by simp [allowed], tr_if_pos (tr_Q4_mem j), Nat.cast_sum]
      rw [pow_zero, Matrix.one_mulVec]
      rw [← tr_P5_image, Finset.sum_image tr_P5_injOn]
      rw [show (Bmat *ᵥ ones5) j
          = ∑ i : Fin 5, (if Disjoint (Q4 j) (P5 i) then (1 : ℤ) else 0) * ones5 i from rfl]
      refine Finset.sum_congr rfl fun i _ => ?_
      have h1 : tr_fillEnd 1 (P5 i) = 1 := by
        rw [tr_fillEnd_one, show allowed 0 = Pprof from by simp [allowed],
          tr_if_pos (tr_P5_mem i)]
      rw [tr_cast_ite_mul, h1, Nat.cast_one, mul_one,
        show ones5 i = (1 : ℤ) from rfl, mul_one]
      exact tr_ite_disjoint_comm (P5 i) (Q4 j)
  | succ k ih =>
    obtain ⟨ihP, ihQ⟩ := ih
    have hallowed1 : allowed (2 * k + 1 + 1) = Pprof := by
      rw [show 2 * k + 1 + 1 = 2 * (k + 1) from by omega]
      exact tr_allowed_double (k + 1)
    have hallowed2 : allowed (2 * k + 2 + 1) = Qprof := by
      rw [show 2 * k + 2 + 1 = 2 * (k + 1) + 1 from by omega]
      exact tr_allowed_double_succ (k + 1)
    have hallowed3 : allowed (2 * k + 2) = Pprof := by
      rw [show 2 * k + 2 = 2 * (k + 1) from by omega]
      exact tr_allowed_double (k + 1)
    have hP1 : ∀ i : Fin 5,
        (tr_fillEnd (2 * (k + 1) + 1) (P5 i) : ℤ) = ((Tmat ^ (k + 1)) *ᵥ ones5) i := by
      intro i
      have key : ((Tmat ^ (k + 1)) *ᵥ ones5) i
          = ∑ j : Fin 4, Bmat j i * (Bmat *ᵥ ((Tmat ^ k) *ᵥ ones5)) j := by
        have h1 : (Amat *ᵥ (Bmat *ᵥ ((Tmat ^ k) *ᵥ ones5))) = (Tmat *ᵥ ((Tmat ^ k) *ᵥ ones5)) :=
          Matrix.mulVec_mulVec _ Amat Bmat
        rw [tr_pow_succ', ← Matrix.mulVec_mulVec, ← h1]
        show (∑ j : Fin 4, Amat i j * (Bmat *ᵥ ((Tmat ^ k) *ᵥ ones5)) j) = _
        exact Finset.sum_congr rfl fun j _ => by rw [tr_Amat_eq_Bmat i j]
      rw [show 2 * (k + 1) + 1 = 2 * k + 1 + 2 from by omega, tr_fillEnd_step, tr_cast_if,
        hallowed1, tr_allowed_double_succ, tr_if_pos (tr_P5_mem i), Nat.cast_sum, key]
      rw [← tr_Q4_image, Finset.sum_image tr_Q4_injOn]
      refine Finset.sum_congr rfl fun j _ => ?_
      rw [tr_cast_ite_mul, show 2 * k + 1 + 1 = 2 * k + 2 from by omega, ihQ j]
      rfl
    refine ⟨hP1, ?_⟩
    intro j
    rw [show (Bmat *ᵥ ((Tmat ^ (k + 1)) *ᵥ ones5)) j
        = ∑ i : Fin 5, (if Disjoint (Q4 j) (P5 i) then (1 : ℤ) else 0)
            * ((Tmat ^ (k + 1)) *ᵥ ones5) i from rfl]
    rw [show 2 * (k + 1) + 2 = 2 * k + 2 + 2 from by omega, tr_fillEnd_step, tr_cast_if,
      hallowed2, hallowed3, tr_if_pos (tr_Q4_mem j), Nat.cast_sum]
    rw [← tr_P5_image, Finset.sum_image tr_P5_injOn]
    refine Finset.sum_congr rfl fun i _ => ?_
    rw [tr_cast_ite_mul, show 2 * k + 2 + 1 = 2 * (k + 1) + 1 from by omega, hP1 i,
      tr_ite_disjoint_comm (P5 i) (Q4 j)]

/-- `tr_Vv` 与矩阵幂表示一致（对 k 归纳）。 -/
theorem tr_Vv_eq (k : ℕ) : tr_Vv k = (Tmat ^ k) *ᵥ ones5 := by
  induction k with
  | zero =>
      rw [pow_zero, Matrix.one_mulVec]
      rfl
  | succ k ih =>
      rw [tr_Vv, ih, Matrix.mulVec_mulVec, tr_pow_succ']

/-- B8：奇列末的填法数 = `Tmat ^ k *ᵥ ones5` 分量。 -/
theorem tr_fillEnd_P5 (k : ℕ) (i : Fin 5) :
    tr_fillEnd (2 * k + 1) (P5 i) = ((Tmat ^ k) *ᵥ ones5) i :=
  (tr_fillEnd_pair k).1 i

/-- B9：偶列末的填法数 = `Bmat *ᵥ (Tmat ^ k *ᵥ ones5)` 分量。 -/
theorem tr_fillEnd_Q4 (k : ℕ) (j : Fin 4) :
    tr_fillEnd (2 * k + 2) (Q4 j) = (Bmat *ᵥ ((Tmat ^ k) *ᵥ ones5)) j :=
  (tr_fillEnd_pair k).2 j

/-- B10：奇宽度计数 = ones5 侧矩阵和。 -/
theorem tr_validCount_odd (k : ℕ) :
    validCount (2 * k + 1) = ∑ i : Fin 5, ((Tmat ^ k) *ᵥ ones5) i := by
  rw [tr_validCount_succ, Nat.cast_sum, tr_allowed_double, ← tr_P5_image,
    Finset.sum_image tr_P5_injOn]
  exact Finset.sum_congr rfl fun i _ => tr_fillEnd_P5 k i

/-- B11：偶宽度计数 = vec0 侧矩阵和。 -/
theorem tr_validCount_even (k : ℕ) :
    validCount (2 * k + 2) = ∑ j : Fin 4, (Bmat *ᵥ ((Tmat ^ k) *ᵥ ones5)) j := by
  rw [show 2 * k + 2 = (2 * k + 1) + 1 from by omega, tr_validCount_succ, Nat.cast_sum,
    show (2 * k + 1) + 1 = 2 * k + 2 from by omega,
    tr_allowed_double_succ, ← tr_Q4_image, Finset.sum_image tr_Q4_injOn]
  exact Finset.sum_congr rfl fun j _ => tr_fillEnd_Q4 k j

-- ============================================================
-- 段 C：代数归约（组装员追加；新引理统一前缀 `asm_`，不与 br_ / tr_ 重名）
-- 数学路线：attempts/guide.md §5（证书 → w/u 递推 → 奇偶归约 → 唯一边界 n = 6）
-- ============================================================

/-!
### C0. 求和 / `*ᵥ` 的 ℤ-线性辅助

不复用 mathlib 分配律名（避免名字漂移），全部自证：`Finset.induction`
（`empty` / `insert` 两情形）+ `Finset.sum_insert` + `ring`。
`asm_mulVec_lin` 经 `Matrix.mulVec`（`dotProduct` 的点态展开）把 `M *ᵥ`
搬到分量上，再套 C0 的求和恒等式。
-/

/-- 常数倍过求和（左系数版）。 -/
theorem asm_mul_sum {ι : Type} (s : Finset ι) (c : ℤ) (g : ι → ℤ) :
    c * ∑ j ∈ s, g j = ∑ j ∈ s, c * g j := by
  classical
  induction s using Finset.induction with
  | empty => simp
  | insert a s ha ih => rw [Finset.sum_insert ha, Finset.sum_insert ha, mul_add, ih]

/-- 两个求和的加法合并。 -/
theorem asm_sum_add {ι : Type} (s : Finset ι) (f g : ι → ℤ) :
    (∑ j ∈ s, f j) + (∑ j ∈ s, g j) = ∑ j ∈ s, (f j + g j) := by
  classical
  induction s using Finset.induction with
  | empty => simp
  | insert a s ha ih =>
      rw [Finset.sum_insert ha, Finset.sum_insert ha, Finset.sum_insert ha]
      calc (f a + ∑ x ∈ s, f x) + (g a + ∑ x ∈ s, g x)
          = (f a + g a) + (∑ x ∈ s, f x + ∑ x ∈ s, g x) := by ring
        _ = (f a + g a) + ∑ x ∈ s, (f x + g x) := by rw [ih]

/-- 两个求和的减法合并。 -/
theorem asm_sum_sub {ι : Type} (s : Finset ι) (f g : ι → ℤ) :
    (∑ j ∈ s, f j) - (∑ j ∈ s, g j) = ∑ j ∈ s, (f j - g j) := by
  classical
  induction s using Finset.induction with
  | empty => simp
  | insert a s ha ih =>
      rw [Finset.sum_insert ha, Finset.sum_insert ha, Finset.sum_insert ha]
      calc (f a + ∑ x ∈ s, f x) - (g a + ∑ x ∈ s, g x)
          = (f a - g a) + (∑ x ∈ s, f x - ∑ x ∈ s, g x) := by ring
        _ = (f a - g a) + ∑ x ∈ s, (f x - g x) := by rw [ih]

/-- ℤ-线性组合 `12 · a - 15 · b + 2 · c` 过求和。 -/
theorem asm_sum_lin {ι : Type} (s : Finset ι) (M : ι → ℤ) (a b c : ι → ℤ) :
    ∑ j ∈ s, M j * (12 * a j - 15 * b j + 2 * c j)
      = 12 * ∑ j ∈ s, M j * a j - 15 * ∑ j ∈ s, M j * b j + 2 * ∑ j ∈ s, M j * c j := by
  rw [asm_mul_sum, asm_mul_sum, asm_mul_sum, asm_sum_sub, asm_sum_add]
  exact Finset.sum_congr rfl fun j _ => by ring

/-!
### C1. w-递推 / u-递推（guide §5.1）

`cert_ones`（interface 段三）给 `k = 0` 的点态递推；`Matrix.mulVec_mulVec`
（向量占第一个显式参数位）+ 裸 `pow_add` 把 `Tmat ^ k *ᵥ` 前后搬，
`asm_mulVec_lin` 处理线性组合，得任意 `k` 的递推。u-递推由 w-递推对
`Bmat *ᵥ` 逐点搬得到（guide §5.1 末段）。
-/

/-- `M *ᵥ` 对点态 ℤ-线性组合的作用（分量形）。 -/
theorem asm_mulVec_lin (M : Matrix (Fin 5) (Fin 5) ℤ) (a b c : Fin 5 → ℤ) (i : Fin 5) :
    (M *ᵥ fun j => 12 * a j - 15 * b j + 2 * c j) i
      = 12 * (M *ᵥ a) i - 15 * (M *ᵥ b) i + 2 * (M *ᵥ c) i := by
  show ∑ j : Fin 5, M i j * (12 * a j - 15 * b j + 2 * c j)
      = 12 * (∑ j : Fin 5, M i j * a j) - 15 * (∑ j : Fin 5, M i j * b j)
        + 2 * (∑ j : Fin 5, M i j * c j)
  exact asm_sum_lin (s := (univ : Finset (Fin 5))) (fun j => M i j) a b c

/-- w-递推（点态）：`w k = Tmat ^ k *ᵥ ones5`。 -/
theorem asm_w_rec (k : ℕ) (i : Fin 5) :
    ((Tmat ^ (k + 3)) *ᵥ ones5) i
      = 12 * ((Tmat ^ (k + 2)) *ᵥ ones5) i - 15 * ((Tmat ^ (k + 1)) *ᵥ ones5) i
        + 2 * ((Tmat ^ k) *ᵥ ones5) i := by
  have hstep : (Tmat ^ (k + 3)) *ᵥ ones5 = (Tmat ^ k) *ᵥ ((Tmat ^ 3) *ᵥ ones5) := by
    rw [pow_add]
    exact (Matrix.mulVec_mulVec _ _ _).symm
  have hcert : (Tmat ^ 3) *ᵥ ones5
      = fun j => 12 * ((Tmat ^ 2) *ᵥ ones5) j - 15 * ((Tmat ^ 1) *ᵥ ones5) j + 2 * ones5 j := by
    funext j
    exact cert_ones j
  have e2 : (Tmat ^ k) *ᵥ ((Tmat ^ 2) *ᵥ ones5) = (Tmat ^ (k + 2)) *ᵥ ones5 := by
    rw [Matrix.mulVec_mulVec, ← pow_add]
  have e1 : (Tmat ^ k) *ᵥ ((Tmat ^ 1) *ᵥ ones5) = (Tmat ^ (k + 1)) *ᵥ ones5 := by
    rw [Matrix.mulVec_mulVec, ← pow_add]
  rw [hstep, hcert, asm_mulVec_lin, e2, e1]

/-- u-递推（点态）：`u k = Bmat *ᵥ (Tmat ^ k *ᵥ ones5)`。 -/
theorem asm_u_rec (k : ℕ) (j : Fin 4) :
    (Bmat *ᵥ ((Tmat ^ (k + 3)) *ᵥ ones5)) j
      = 12 * (Bmat *ᵥ ((Tmat ^ (k + 2)) *ᵥ ones5)) j
        - 15 * (Bmat *ᵥ ((Tmat ^ (k + 1)) *ᵥ ones5)) j + 2 * (Bmat *ᵥ ((Tmat ^ k) *ᵥ ones5)) j := by
  show ∑ i : Fin 5, Bmat j i * (((Tmat ^ (k + 3)) *ᵥ ones5) i)
      = 12 * (∑ i : Fin 5, Bmat j i * (((Tmat ^ (k + 2)) *ᵥ ones5) i))
        - 15 * (∑ i : Fin 5, Bmat j i * (((Tmat ^ (k + 1)) *ᵥ ones5) i))
        + 2 * (∑ i : Fin 5, Bmat j i * (((Tmat ^ k) *ᵥ ones5) i))
  rw [asm_mul_sum, asm_mul_sum, asm_mul_sum, asm_sum_sub, asm_sum_add]
  exact Finset.sum_congr rfl fun i _ => by rw [asm_w_rec k i]; ring

/-- w-递推（求和形，奇侧用）。 -/
theorem asm_w_rec_sum (k : ℕ) :
    (∑ i : Fin 5, ((Tmat ^ (k + 3)) *ᵥ ones5) i)
      = 12 * (∑ i : Fin 5, ((Tmat ^ (k + 2)) *ᵥ ones5) i)
        - 15 * (∑ i : Fin 5, ((Tmat ^ (k + 1)) *ᵥ ones5) i) + 2 * (∑ i : Fin 5, ((Tmat ^ k) *ᵥ ones5) i) := by
  rw [asm_mul_sum, asm_mul_sum, asm_mul_sum, asm_sum_sub, asm_sum_add]
  exact Finset.sum_congr rfl fun i _ => asm_w_rec k i

/-- u-递推（求和形，偶侧用）。 -/
theorem asm_u_rec_sum (k : ℕ) :
    (∑ j : Fin 4, (Bmat *ᵥ ((Tmat ^ (k + 3)) *ᵥ ones5)) j)
      = 12 * (∑ j : Fin 4, (Bmat *ᵥ ((Tmat ^ (k + 2)) *ᵥ ones5)) j)
        - 15 * (∑ j : Fin 4, (Bmat *ᵥ ((Tmat ^ (k + 1)) *ᵥ ones5)) j)
        + 2 * (∑ j : Fin 4, (Bmat *ᵥ ((Tmat ^ k) *ᵥ ones5)) j) := by
  rw [asm_mul_sum, asm_mul_sum, asm_mul_sum, asm_sum_sub, asm_sum_add]
  exact Finset.sum_congr rfl fun j _ => asm_u_rec k j

/-!
### C2. 主定理（陈述与 `formalized/statement.lean` 冻结件逐字一致）

按 guide §5.2–5.3 组装：
- 奇 `n = 2m + 1`（`6 ≤ n` ⟹ `m ≥ 3`）：走 w-递推，无边界情形；
- 偶 `n = 2m`（`n ≠ 6` ⟹ `m ≥ 4`）：走 u-递推；
- 唯一边界 `n = 6`：`z 6 / z 4 / z 2` 走矩阵和、`z 0 = 1` 走 B6，`decide` 闭。

ℕ/ℤ 口径：`tr_validCount_odd` / `tr_validCount_even` / `tr_validCount_zero`
的右端是 ℤ 矩阵和，故每个分支先把要用的等式搬到 ℤ（`have key : (… : ℤ) = …`），
分支内用 `omega` 收口；最后 `omega` 借 ℤ 假设回证 ℕ 目标（本机 kernel 实测
`omega` 可桥接 `Nat.cast` 两侧，见 `tmp-probe/comb12-asm-probe.lean` T2）。
`n - 4` / `n - 2` / `n - 6` 是 ℕ 截断减法，用 `omega` 翻译成 `2 * (m - k) + c` 形。
-/

theorem c_z_rec : ∀ n : ℕ, 6 ≤ n → z n + 15 * z (n - 4) = 12 * z (n - 2) + 2 * z (n - 6) := by
  intro n hn
  rcases eq_or_ne n 6 with rfl | hne
  · -- 唯一边界 n = 6
    have h4 : (6 - 4 : ℕ) = 2 := by decide
    have h2 : (6 - 2 : ℕ) = 4 := by decide
    have h0 : (6 - 6 : ℕ) = 0 := by decide
    have key : (z 6 : ℤ) + 15 * (z 2 : ℤ) = 12 * (z 4 : ℤ) + 2 * (z 0 : ℤ) := by
      have h1 : (z 6 : ℤ) = ∑ j : Fin 4, (Bmat *ᵥ ((Tmat ^ 2) *ᵥ ones5)) j := by
        have h := tr_validCount_even 2
        rw [← br_z_eq_validCount (2 * 2 + 2)] at h
        exact h
      have h2' : (z 2 : ℤ) = ∑ j : Fin 4, (Bmat *ᵥ ((Tmat ^ 0) *ᵥ ones5)) j := by
        have h := tr_validCount_even 0
        rw [← br_z_eq_validCount (2 * 0 + 2)] at h
        exact h
      have h3 : (z 4 : ℤ) = ∑ j : Fin 4, (Bmat *ᵥ ((Tmat ^ 1) *ᵥ ones5)) j := by
        have h := tr_validCount_even 1
        rw [← br_z_eq_validCount (2 * 1 + 2)] at h
        exact h
      have h0z : (z 0 : ℤ) = 1 := by
        have h := tr_validCount_zero
        rw [← br_z_eq_validCount 0] at h
        exact_mod_cast h
      have vsum : (∑ j : Fin 4, (Bmat *ᵥ ((Tmat ^ 2) *ᵥ ones5)) j)
          + 15 * (∑ j : Fin 4, (Bmat *ᵥ ((Tmat ^ 0) *ᵥ ones5)) j)
          = 12 * (∑ j : Fin 4, (Bmat *ᵥ ((Tmat ^ 1) *ᵥ ones5)) j) + 2 * 1 := by
        decide
      omega
    rw [h4, h2, h0]
    omega
  · -- n ≥ 7：奇偶分岔（guide §5.2）
    have hpar : ∃ m, n = 2 * m ∨ n = 2 * m + 1 := by
      refine ⟨n / 2, ?_⟩
      omega
    obtain ⟨m, hw⟩ := hpar
    rcases hw with rfl | rfl
    · -- 偶 n = 2 * m，hne : 2 * m ≠ 6 且 6 ≤ 2 * m ⟹ m ≥ 4
      have hm : 4 ≤ m := by omega
      have e4 : 2 * m - 4 = 2 * (m - 3) + 2 := by omega
      have e2 : 2 * m - 2 = 2 * (m - 2) + 2 := by omega
      have e6 : 2 * m - 6 = 2 * (m - 4) + 2 := by omega
      have e0 : 2 * m = 2 * (m - 1) + 2 := by omega
      have key : (z (2 * (m - 1) + 2) : ℤ) + 15 * (z (2 * (m - 3) + 2) : ℤ)
          = 12 * (z (2 * (m - 2) + 2) : ℤ) + 2 * (z (2 * (m - 4) + 2) : ℤ) := by
        have h1 : (z (2 * (m - 1) + 2) : ℤ)
            = ∑ j : Fin 4, (Bmat *ᵥ ((Tmat ^ (m - 1)) *ᵥ ones5)) j := by
          have h := tr_validCount_even (m - 1)
          rw [← br_z_eq_validCount (2 * (m - 1) + 2)] at h
          exact h
        have h2 : (z (2 * (m - 3) + 2) : ℤ)
            = ∑ j : Fin 4, (Bmat *ᵥ ((Tmat ^ (m - 3)) *ᵥ ones5)) j := by
          have h := tr_validCount_even (m - 3)
          rw [← br_z_eq_validCount (2 * (m - 3) + 2)] at h
          exact h
        have h3 : (z (2 * (m - 2) + 2) : ℤ)
            = ∑ j : Fin 4, (Bmat *ᵥ ((Tmat ^ (m - 2)) *ᵥ ones5)) j := by
          have h := tr_validCount_even (m - 2)
          rw [← br_z_eq_validCount (2 * (m - 2) + 2)] at h
          exact h
        have h4 : (z (2 * (m - 4) + 2) : ℤ)
            = ∑ j : Fin 4, (Bmat *ᵥ ((Tmat ^ (m - 4)) *ᵥ ones5)) j := by
          have h := tr_validCount_even (m - 4)
          rw [← br_z_eq_validCount (2 * (m - 4) + 2)] at h
          exact h
        have hrec := asm_u_rec_sum (m - 4)
        rw [show (m - 4) + 3 = m - 1 from by omega, show (m - 4) + 2 = m - 2 from by omega,
          show (m - 4) + 1 = m - 3 from by omega] at hrec
        omega
      rw [e4, e2, e6, e0]
      omega
    · -- 奇 n = 2 * m + 1，6 ≤ 2 * m + 1 ⟹ m ≥ 3
      have hm : 3 ≤ m := by omega
      have e4 : 2 * m + 1 - 4 = 2 * (m - 2) + 1 := by omega
      have e2 : 2 * m + 1 - 2 = 2 * (m - 1) + 1 := by omega
      have e6 : 2 * m + 1 - 6 = 2 * (m - 3) + 1 := by omega
      have key : (z (2 * m + 1) : ℤ) + 15 * (z (2 * (m - 2) + 1) : ℤ)
          = 12 * (z (2 * (m - 1) + 1) : ℤ) + 2 * (z (2 * (m - 3) + 1) : ℤ) := by
        have h1 : (z (2 * m + 1) : ℤ) = ∑ i : Fin 5, ((Tmat ^ m) *ᵥ ones5) i := by
          have h := tr_validCount_odd m
          rw [← br_z_eq_validCount (2 * m + 1)] at h
          exact h
        have h2 : (z (2 * (m - 2) + 1) : ℤ) = ∑ i : Fin 5, ((Tmat ^ (m - 2)) *ᵥ ones5) i := by
          have h := tr_validCount_odd (m - 2)
          rw [← br_z_eq_validCount (2 * (m - 2) + 1)] at h
          exact h
        have h3 : (z (2 * (m - 1) + 1) : ℤ) = ∑ i : Fin 5, ((Tmat ^ (m - 1)) *ᵥ ones5) i := by
          have h := tr_validCount_odd (m - 1)
          rw [← br_z_eq_validCount (2 * (m - 1) + 1)] at h
          exact h
        have h4 : (z (2 * (m - 3) + 1) : ℤ) = ∑ i : Fin 5, ((Tmat ^ (m - 3)) *ᵥ ones5) i := by
          have h := tr_validCount_odd (m - 3)
          rw [← br_z_eq_validCount (2 * (m - 3) + 1)] at h
          exact h
        have hrec := asm_w_rec_sum (m - 3)
        rw [show (m - 3) + 3 = m from by omega, show (m - 3) + 2 = m - 1 from by omega,
          show (m - 3) + 1 = m - 2 from by omega] at hrec
        omega
      rw [e4, e2, e6]
      omega

#print axioms c_z_rec
\end{lstlisting}

\section{Verification evidence}\label{sec:verify}

\begin{itemize}
\item Toolchain: Lean \texttt{v4.34.0} (\texttt{leanprover/lean4:v4.34.0}),
mathlib4 \texttt{v4.34.0}, revision \texttt{5ed29652}. Reproduction: with
that pin, a strict compile (0 \texttt{sorry}, 0 warning) of
\texttt{proofs/02-grid3n-colmid-}\allowbreak\texttt{indep-proved.lean}
exits with code $0$ and
empty output apart from the interface roster's \texttt{\#check} echoes.
\item Statement integrity: the frozen statement file (61 content lines) has sha256
\begin{center}
\texttt{b726f0d8e01672b39a93ce0c256f5036014}\allowbreak
\texttt{bc80f7f0b6a5fcb484272d805c884},
\end{center}
the frozen interface file (8042 bytes) has sha256
\begin{center}
\texttt{66c96c3e4051ce2be5bd2c09383913f71678}\allowbreak
\texttt{e3133b0e67cff4fdc73dbed3d6e9},
\end{center}
and the final proof file (1188 lines) has sha256
\begin{center}
\texttt{9a8378352c7afe91ba83a519053e130c53}\allowbreak
\texttt{bd6fd7d08282c48669289682f751ac}.
\end{center}
The audit department's byte-exactness check reports: the first 8042 bytes of
the final file equal the interface file byte for byte (this embeds the
frozen statement segment verbatim, since the interface file opens with it),
and the theorem statement line in the final file is byte-equal to line 48 of
the frozen statement file. The declaration \texttt{c\_z\_rec} occurs exactly
twice in the final file (the statement and the \texttt{\#print axioms}
command) --- no second, altered statement is smuggled in.
\item Kernel check: compiles with $0$ \texttt{sorry} (whole-file scan
including comments, also covering \texttt{admit}: zero hits);
\texttt{\#print axioms} output, verbatim (kernel re-run of 2026-09-30,
exit code 0, kept in \texttt{audit/axioms.out}):
\begin{lstlisting}[basicstyle=\ttfamily\footnotesize,frame=none]
'c_z_rec' depends on axioms: [propext, Classical.choice, Quot.sound]
\end{lstlisting}
This equals the whitelist \{propext, Quot.sound, Classical.choice\} exactly;
in particular no \texttt{sorryAx} and no \texttt{Lean.ofReduceBool} (the
\texttt{native\_decide} axiom) appear. The single textual occurrence of
\texttt{native\_decide} in the final file is inside the interface's warning
docstring, which \emph{bans} it; the axiom output is the authoritative
criterion.
\item Grading by the audit department: the single proven declaration is
\emph{complete proof}; the lowest grade reached for the case = \emph{complete
proof} (final adjudication \texttt{final-audit.md}: byte-exactness checks,
banned-word scan, non-vacuous-lemma spot checks on the bridge and transfer
lanes, axiom set check --- all passed). No weakening, no added hypotheses,
no unproven residue.
\item Audit chain (originals in this deposit under \texttt{audit/} unless
noted): the final adjudication, the axiom re-print, this note's claim check
\texttt{claim-check.md} --- added to \texttt{audit/} after the audit
department returns its adjudication, and a prerequisite for publication.
\end{itemize}

\section{Novelty evidence}\label{sec:novelty}

Adjudication copied from the pipeline's exclusivity review (C9 gate), which
ran before the problem was accepted: verdict \textbf{no occupying record
found} --- at the object level (the counting sequence $z$ of this specific
defective strip is unregistered) and at the proposition level (no record
states this recurrence for this object). Value tier \textbf{new sequence,
lower tier}: the ``new sequence'' half holds ($z$ itself is unregistered in
OEIS across five query forms, and the defective-strip object is unnamed in
the literature layer), while the ``new recurrence'' half does \emph{not}
hold and is not claimed: the coefficient signature $(12,-15,2)$ of the
parity subsequences' order-3 recurrences is already indexed in the OEIS
linear-recurrence index, and the even-indexed subsequence
$z(2), z(4), \dots$ coincides term-by-term (over their overlap window) with
the registered sequence A122011 \cite{a122011} --- a boundary that must be
stated whenever that subsequence is mentioned. A122011's own record is a
$3\times3$ matrix-power exercise with no graph-theoretic or
independent-set content, so the object is different even where the values
agree; the odd-indexed subsequence has no registered match. The nearest
literature neighbours are delimited: the Merrifield--Simmons line
\cite{trinks2021} treats the quantity $\sigma(G - S)$ for vertex-deleted
graphs as a tool in sign-relational results (not this counting sequence),
the De~Ita et al.\ branching work \cite{deita2024} counts independent sets
of \emph{intact} grid graphs by vertex elimination (no deleted-column
family), and the independence-complex work \cite{jonsson2007} concerns
square-grid complexes (no counting sequence). Summary by channel; the full
query log, including every zero-hit query, is in this deposit at
\texttt{audit/c9-record.md} (verbatim, byte-for-byte).

\begin{center}
\resizebox{\textwidth}{!}{%
\begin{tabular}{lll}
\hline
channel & queries & result \\
\hline
OEIS, numeric & 4 prefix-window strings (main; drop-first-1; drop-first-2;
odd-indexed subsequence) & all zero hits; positive-control anchor hit
(A051736, the intact-strip sequence) \\
OEIS, numeric (even subsequence) & 1 string & hit A122011: same values,
different object (matrix-power example); recurrence signature already
indexed there \\
OEIS, keyword & 6 keyword strings (defective/periodic grid, independent
sets, deleted vertices) & no this-object hit \\
arXiv API & 11 strings (object-name $\times$ property words, plus variants)
& all returned; Top-3 candidates read at full-text level, none this object \\
Crossref & 4 strings & one high-relevance record (De~Ita et al.\ 2024,
intact grids); not this object \\
Mathematics Stack Exchange & 2 site-scoped strings & no this-object hit \\
local formalisation ecosystem & 6 GitHub strings + 4 mathlib \texttt{rg}
sweeps & counting object absent; no same-form formalisation \\
residual exposure & --- & the literature-layer zero-hit verdict is limited
to arXiv + Crossref + MSE; OpenAlex / zbMATH / general web search were not
reachable \\
\hline
\end{tabular}%
}
\end{center}

The channels are named so a later reader can re-run them; if any channel
later returns a same-form record, the tier above weakens accordingly, and
this deposit will be amended by a later version rather than by editing this
record.

\section{Scope and limitations}\label{sec:limits}

Grade and boundary, stated as the audit states them. (a) The kernel layer
proves Theorem~\ref{thm:rec} and nothing else. (b) The grading is
\emph{complete proof} for that theorem only; this note makes no claim about
recurrence minimality (that no shorter-order recurrence holds is neither
proved nor searched), about closed forms, or about asymptotics. (c) The
identification of \texttt{z} with the independent-set count of the
defective strip is definitional in the formalisation (\texttt{z} is defined
as that count), and the bridge lemma is kernel-proved; no further
combinatorial conjecture layer is carried. (d) The novelty verdict's
literature-layer half is limited to three sources (arXiv, Crossref, MSE);
OpenAlex, zbMATH and general web search were not reachable in the review.
(e) The initial values stated in the frozen statement file were verified
there by \texttt{native\_decide}, which is a data check outside the axiom
whitelist and deliberately excluded from the deposited proof; the deposited
proof derives everything it needs through the transfer-matrix route.
(f) The gate-four sign-off of the underlying task was exercised by the
author's instruction of 2026-09-30 to run both the paper lane and this
claim lane with GitHub outbound authorisation, while the task report file
header retains its pre-sign-off wording unchanged; both calibers are
recorded. (g) This note carries no literature review and is a deposit
record, not a survey; the narrative paper is prepared separately under the
pipeline's paper lane and will reference this deposit.

\section*{Data availability}

The deposit package for this claim contains the claim note (LaTeX source and
PDF), the Lean sources (\texttt{proofs/}: frozen statement file, frozen
interface file, final proof file, \texttt{lean-toolchain}), the verification
and audit logs, and the full C9 query log with zero-hit entries
(\texttt{audit/}). Public mirror: the claim lane's repository
\url{https://github.com/Aurora0134/ai4math-claims}, directory
\texttt{grid3n-colmid-indep/} (pushed under the author's instruction of
2026-09-30; the commit and timestamp are recorded in the source
repository's \texttt{claims/grid3n-colmid-indep/card.md}). Zenodo: the DOI
is reserved at upload by the author and is recorded in
\texttt{metadata/README.md} of this package once created; see
\texttt{UPLOAD.md} in the source repository. Claim note under CC BY 4.0,
code under MIT.

\section*{Use of generative AI}

(a) AI performed: candidate generation and the exclusivity review,
autoformalisation of the definitions and statement, the numerical
cross-check of the proof route (three independent counting implementations
over $n = 0..20$), proof search and proof-script writing (eight worker
agents in a workflow: blueprint, blueprint review, bridge, transfer,
assembly, final audit, report, reader audit), and the assembly of this note
including its verbatim listings. (b) Model, node and budget: the proving
phase ran on the session node of the pipeline host (the wire designation
\texttt{a6api-main/glm-5.3-flash} is recorded in the claim-lane card's node
snapshot,
\texttt{claims/grid3n-colmid-indep/card.md}; the source task ledger records
the same session's node as its session node), with no cross-node switching,
consuming $8$ agent-run
samples of a permitted $16$, $0$ llm-calls of a permitted $8$, and $10$
strict kernel compiles of a permitted $36$ (plus compile-debug probes,
separately ledgered), copied verbatim from the task ledger. This note
consumed $0$ pipeline LLM calls against a cap of $4$; drafting was
mechanical assembly. (c) Final adjudication: all statements and proofs
reported here were checked by the Lean 4 kernel: the proof compiles with no
\texttt{sorry}, and \texttt{\#print axioms} reports only \{propext,
Quot.sound, Classical.choice\}. (d) The AI system is not an author; the
human author reviewed the evidence and is responsible for all content.

\section*{Author contributions}

CRediT-style: the human author adjudicated topic selection, statement
semantics and the deposit decision; the AI4Math pipeline (an LLM) generated
the candidate propositions, the proof searches and the text drafts. The AI
system is not listed as an author.

\begin{thebibliography}{9}
% Each entry was verified live on 2026-09-30 (lean4, mathlib: Crossref
% title/page metadata; a122011, a051736: live oeis.org fetches; trinks2021,
% jonsson2007: arXiv abs pages; deita2024: Crossref DOI record). No entry is
% quoted from memory.
\bibitem{lean4}
L.~de Moura and S.~Ullrich.
\newblock The Lean 4 theorem prover and programming language.
\newblock \emph{Automated Deduction -- CADE-28}, LNCS, pp.~625--635, 2021.
\url{https://doi.org/10.1007/978-3-030-79876-5_37}

\bibitem{mathlib}
The mathlib Community.
\newblock The Lean mathematical library.
\newblock \emph{CPP 2020}, pp.~367--381, 2020.
\url{https://doi.org/10.1145/3372885.3373824}

\bibitem{a122011}
OEIS Foundation.
\newblock Sequence A122011 (generating function $x^2(1+x)/(1-12x+15x^2-2x^3)$;
the $3\times3$ matrix-power example of A.~M.~Robert, \emph{Linear Algebra,
Examples and Applications}, World Scientific, 2005, p.~58; linear-recurrence
index, signature $(12,-15,2)$).
\newblock \emph{The On-Line Encyclopedia of Integer Sequences}, entry
A122011.
\url{https://oeis.org/A122011}

\bibitem{a051736}
OEIS Foundation.
\newblock Sequence A051736 (number of independent sets of the $3\times n$
grid strip; positive-control anchor of the review).
\newblock \emph{The On-Line Encyclopedia of Integer Sequences}, entry
A051736.
\url{https://oeis.org/A051736}

\bibitem{trinks2021}
M.~Trinks.
\newblock The Merrifield--Simmons conjecture holds for bipartite graphs.
\newblock \emph{arXiv:1006.4253}, 2021 (v2).
\url{https://arxiv.org/abs/1006.4253}

\bibitem{deita2024}
G.~De~Ita, M.~Bello-L\'opez, and J.~Marcial-Romero.
\newblock Counting rules for computing the number of independent sets of a
grid graph.
\newblock \emph{Mathematics} 12(6):922, 2024.
\url{https://doi.org/10.3390/math12060922}

\bibitem{jonsson2007}
J.~Jonsson.
\newblock On the independence complex of square grids.
\newblock \emph{arXiv:math/0701890}, 2007 (v2).
\url{https://arxiv.org/abs/math/0701890}

\end{thebibliography}

\end{document}
